% ============================================================
%  VitaChain — Rapport de Projet de Fin d'Études
%  Compilé avec XeLaTeX ou pdfLaTeX sur Overleaf
% ============================================================
\documentclass[12pt,a4paper]{report}

% ── Encodage & langue ─────────────────────────────────────
\usepackage[utf8]{inputenc}
\usepackage[T1]{fontenc}
\usepackage[french]{babel}

% ── Mise en page ──────────────────────────────────────────
\usepackage[top=2.5cm,bottom=2.5cm,left=3cm,right=2.5cm]{geometry}
\usepackage{setspace}
\onehalfspacing

% ── Typographie ───────────────────────────────────────────
\usepackage{lmodern}
\usepackage{microtype}

% ── Maths & symboles ──────────────────────────────────────
\usepackage{amssymb}   % \checkmark, \square

% ── Diagrammes TikZ & Gantt ───────────────────────────────
\usepackage{tikz}
\usetikzlibrary{shapes.geometric,arrows.meta,positioning,calc,patterns}
\usepackage{pgfgantt}
\usepackage{pgf-pie}

% ── Couleurs ──────────────────────────────────────────────
\usepackage[table]{xcolor}
\definecolor{vitaBlue}{RGB}{37,99,235}
\definecolor{vitaGreen}{RGB}{22,163,74}
\definecolor{codeGray}{RGB}{248,248,250}
\definecolor{codeBorder}{RGB}{200,200,210}
\definecolor{rowHead}{RGB}{219,234,254}

% ── Hyperliens (bordure rouge comme dans le modèle) ───────
\usepackage[colorlinks=false,
            pdfborder={0 0 1},
            linkbordercolor={1 0 0},
            urlbordercolor={0 0 1},
            pdftitle={VitaChain -- Rapport PFE},
            pdfauthor={Yasser}]{hyperref}

% ── Figures ───────────────────────────────────────────────
\usepackage{graphicx}
\usepackage{float}
\usepackage{caption}
\captionsetup{font=small,labelfont=bf,margin=12pt}

% ── Tableaux ──────────────────────────────────────────────
\usepackage{booktabs}
\usepackage{tabularx}
\usepackage{longtable}
\usepackage{multirow}
\usepackage{array}
\renewcommand{\arraystretch}{1.35}

% ── Code source ───────────────────────────────────────────
\usepackage{listings}
\lstdefinestyle{sql}{
  language        = SQL,
  basicstyle      = \ttfamily\small,
  backgroundcolor = \color{codeGray},
  frame           = single,
  rulecolor       = \color{codeBorder},
  breaklines      = true,
  keepspaces      = true,
  showstringspaces= false,
  tabsize         = 2,
  numbers         = left,
  numberstyle     = \tiny\color{gray},
  numbersep       = 8pt,
  keywordstyle    = \color{vitaBlue}\bfseries,
  stringstyle     = \color{vitaGreen},
  commentstyle    = \color{gray}\itshape,
  morekeywords    = {RETURNS,LANGUAGE,SECURITY,DEFINER,PERFORM,RAISE,
                     EXCEPTION,NOTIFY,DECLARE,BEGIN,END,UUID,
                     TIMESTAMPTZ,NUMERIC,JSONB,TEXT,BOOLEAN},
}
\lstdefinestyle{json}{
  basicstyle      = \ttfamily\small,
  backgroundcolor = \color{codeGray},
  frame           = single,
  rulecolor       = \color{codeBorder},
  breaklines      = true,
  showstringspaces= false,
  string          = [s]{"}{"},
  stringstyle     = \color{vitaGreen},
  comment         = [l]{//},
  commentstyle    = \color{gray}\itshape,
}
\lstdefinestyle{bash}{
  basicstyle      = \ttfamily\small,
  backgroundcolor = \color{codeGray},
  frame           = single,
  rulecolor       = \color{codeBorder},
  breaklines      = true,
  showstringspaces= false,
  commentstyle    = \color{gray}\itshape,
}

% ── En-têtes ──────────────────────────────────────────────
\usepackage{fancyhdr}
\pagestyle{fancy}
\fancyhf{}
\fancyhead[L]{\small\itshape\leftmark}
\fancyhead[R]{\small VitaChain --- PFE 2026}
\fancyfoot[C]{\thepage}
\renewcommand{\headrulewidth}{0.4pt}
\fancypagestyle{plain}{
  \fancyhf{}
  \fancyfoot[C]{\thepage}
  \renewcommand{\headrulewidth}{0pt}
}

% ── Listes ────────────────────────────────────────────────
\usepackage{enumitem}
\setlist[itemize]{itemsep=2pt,topsep=3pt}
\setlist[enumerate]{itemsep=2pt,topsep=3pt}

% ── Commande placeholder figure ───────────────────────────
% Utilisation : \fig{nom_fichier}{Légende}{label}
% Mettre les captures dans un dossier "figures/" sur Overleaf
% Remplacer le bloc gris par \includegraphics dès que l'image est prête
\newcommand{\fig}[3]{%
  \begin{figure}[H]
    \centering
    \IfFileExists{figures/#1.png}{%
      \includegraphics[width=0.93\textwidth]{figures/#1.png}%
    }{%
      \IfFileExists{figures/#1.jpg}{%
        \includegraphics[width=0.93\textwidth]{figures/#1.jpg}%
      }{%
        \fcolorbox{codeBorder}{codeGray}{%
          \begin{minipage}[c][4cm][c]{0.90\textwidth}
            \centering
            \textcolor{gray}{\Large$\square$}\\[6pt]
            \textcolor{gray}{\small\itshape
              À insérer : \texttt{figures/#1.png}}
          \end{minipage}}%
      }%
    }%
    \caption{#2}
    \label{fig:#3}
  \end{figure}}

% ── Macro commodités ──────────────────────────────────────
\newcommand{\code}[1]{\texttt{\small #1}}
\newcommand{\api}[1]{\texttt{\small\color{vitaBlue}#1}}

% ── Profondeur TOC ────────────────────────────────────────
\setcounter{tocdepth}{3}
\setcounter{secnumdepth}{3}

% =============================================================
\begin{document}

% ── PAGE DE GARDE ─────────────────────────────────────────
\begin{titlepage}
  \centering

  % ── Filet supérieur double ──
  \noindent\textcolor{vitaBlue}{\rule{\linewidth}{4pt}}\\[-2pt]
  \noindent\textcolor{vitaGreen}{\rule{\linewidth}{2pt}}

  \vspace{0.7cm}

  % ── Établissement ──
  \begin{minipage}{0.48\linewidth}
    \flushleft
    \small\bfseries Faculté des sciences Ben M'Sik\\[3pt]
    \small Département Informatique\\[2pt]
    \small Filière~: Filière Devloppement Informatique et Adminstration Reseaux et Systemes
  \end{minipage}\hfill
  \begin{minipage}{0.48\linewidth}
    \flushright
    % \includegraphics[height=1.6cm]{figures/logo_ecole.png}
    \small\bfseries Année universitaire\\[3pt]
    \small 2025 --- 2026
  \end{minipage}

  \vspace{1.2cm}

  % ── Logo VitaChain ──
  \IfFileExists{figures/vitachain_logo.png}{%
    
\includegraphics[height=3.8cm]{figures/vitachain_logo.png}%
  }{%
    {\Huge\bfseries\color{vitaBlue} VitaChain}%
  }

  \vspace{1.0cm}

  % ── Filet médian ──
  \noindent\textcolor{vitaBlue}{\rule{\linewidth}{1.5pt}}

  \vspace{0.7cm}

  {\large\bfseries\color{vitaBlue} Rapport de Projet de Fin d'Études}

  \vspace{0.8cm}

  {\LARGE\bfseries VitaChain}\\[0.5cm]
  {\large Plateforme numérique intégrée pour l'agriculture intelligente}\\[0.3cm]
  {\large et la chaîne d'approvisionnement alimentaire}

  \vspace{0.5cm}
  {\small\color{gray}
    FastAPI~$\cdot$~Next.js~15~$\cdot$~PostgreSQL~17~$\cdot$~Supabase~$\cdot$~ESP32}

  \vspace{0.7cm}

  % ── Modules (4 étiquettes colorées) ──
  
  % ── Filet médian ──
  \noindent\textcolor{vitaBlue}{\rule{\linewidth}{1.5pt}}

  \vfill

  % ── Étudiant / Encadrant ──
  \begin{minipage}{0.46\linewidth}
    \flushleft
    \textbf{Réalisé par~:}\\[6pt]
    {\large Badre Saad}\\[3pt]
    {\large Kodoussi Mohammed}\\[3pt]
    {\large El Karmi Yasser}\\[3pt]
    \small 2026~$\cdot$~Filière Devloppement Informatique et Adminstration Reseaux et Systemes
  \end{minipage}\hfill
  \begin{minipage}{0.46\linewidth}
    \flushright
    \textbf{Encadré par~:}\\[6pt]
    {\large Ouchra Hafssa}\\[3pt]
    {\large Achtaich Khadija}\\[3pt]
    {\large Ait Daoud Mohammed}\\[3pt]
    \small [Grade~$\cdot$~Établissement]
  \end{minipage}

  \vspace{0.9cm}
  {\small\color{gray} Soutenu le 19/06/2026}

  \vspace{0.5cm}

  % ── Filet inférieur double ──
  \noindent\textcolor{vitaGreen}{\rule{\linewidth}{2pt}}\\[-2pt]
  \noindent\textcolor{vitaBlue}{\rule{\linewidth}{4pt}}

\end{titlepage}

% ── PAGES LIMINAIRES (numérotation romaine) ───────────────
\pagenumbering{roman}
\setcounter{page}{2}

% ── DÉDICACES ─────────────────────────────────────────────
\chapter*{Dédicaces}
\addcontentsline{toc}{chapter}{Dédicaces}

À mes chers parents, pour leur soutien indéfectible, leur amour et leurs sacrifices
tout au long de mon parcours académique. Aucun mot ne saurait exprimer ma profonde
gratitude et mon amour sincère.

À mes frères et sœurs, pour leur encouragement constant et leur présence à mes côtés.

À mes enseignants ainsi qu'à toutes les personnes qui m'ont encouragé durant ce projet.

À tous mes amis et collègues, qui ont partagé avec moi les moments de joie et de
difficulté.

À tous ceux qui me sont chers\ldots

\newpage

% ── REMERCIEMENTS ─────────────────────────────────────────
\chapter*{Remerciements}
\addcontentsline{toc}{chapter}{Remerciements}

Au terme de ce travail, il nous est agréable d'adresser nos sincères
remerciements à toutes les personnes qui ont contribué, de près ou de
loin, à sa réalisation.

Nous tenons à exprimer notre profonde gratitude à notre encadrante,
\textbf{Mme Ouchra Hafsa}, pour ses précieux conseils, sa disponibilité
et ses orientations tout au long de ce projet.

Nous remercions également l'ensemble du corps enseignant de la
\textbf{Faculté des Sciences Ben M'Sik} pour la qualité de la formation
dispensée.

Enfin, nous remercions nos familles et nos amis pour leur soutien moral
tout au long de notre parcours.

\newpage

\tableofcontents\newpage
\listoffigures\newpage
\listoftables\newpage

% ── CORPS DU RAPPORT (numérotation arabe) ─────────────────
\pagenumbering{arabic}
\setcounter{page}{1}

% =============================================================
% CHAPITRE 1 — CAHIER DES CHARGES
% =============================================================
\chapter{Cahier des charges}

\section{Introduction}

Dans un contexte mondial marqué par des défis croissants liés à la sécurité alimentaire,
à la durabilité agricole et à la gestion des ressources naturelles, les technologies
numériques et l'Internet des objets (IoT) offrent des opportunités sans précédent pour
transformer les pratiques du secteur agroalimentaire. Le présent rapport documente la
conception, l'analyse et la réalisation de \textbf{VitaChain}, une plateforme numérique
intégrée dédiée à l'agriculture intelligente et à la chaîne d'approvisionnement
alimentaire au Maroc.

Ce premier chapitre établit le cadre général du projet à travers une présentation du
contexte, des problématiques identifiées, des objectifs visés et des besoins fonctionnels
et non fonctionnels auxquels la solution doit répondre.

\section{Contexte du projet}

Le secteur agricole marocain représente une composante stratégique de l'économie nationale.
Cependant, les agriculteurs marocains, en particulier les petits et moyens exploitants,
font face à des difficultés structurelles : manque d'accès aux données agronomiques en
temps réel, absence de canaux directs vers les acheteurs professionnels, et gaspillage
alimentaire important dans la restauration et la grande distribution.

VitaChain s'inscrit dans le cadre des initiatives de transformation numérique du secteur
agricole. Le projet vise à proposer un écosystème numérique unifié combinant :

\begin{itemize}
  \item \textbf{La surveillance IoT des parcelles agricoles} via des capteurs de sol
        connectés (module Katara) ;
  \item \textbf{Une place de marché agricole directe} entre agriculteurs et restaurateurs
        (module FarMarket) ;
  \item \textbf{Une solution anti-gaspillage alimentaire} permettant aux restaurants de
        redistribuer leurs surplus aux citoyens (module SecondServe) ;
  \item \textbf{Un panneau d'administration centralisé} pour la gouvernance de la
        plateforme.
\end{itemize}

Le projet est réalisé dans le cadre d'un Projet de Fin d'Études (PFE) en ingénierie
logicielle, avec une architecture orientée production déployée sur un VPS Linux avec
PostgreSQL, FastAPI, Next.js et Supabase.

\section{Présentation du problème}

Plusieurs problèmes structurels freinent le développement du secteur agroalimentaire
marocain.

\paragraph{Manque d'informations agronomiques en temps réel.}
Les agriculteurs prennent leurs décisions d'irrigation et de fertilisation sur la base
d'observations empiriques, sans disposer de données précises sur l'état du sol (humidité,
température, pH, conductivité). Cela entraîne une surconsommation d'eau et d'intrants,
réduisant les rendements et augmentant les coûts.

\paragraph{Désintermédiation insuffisante dans la chaîne alimentaire.}
La chaîne de distribution agricole au Maroc repose sur des intermédiaires traditionnels
(grossistes, souk), qui captent une part importante de la valeur ajoutée au détriment
des agriculteurs. Les restaurateurs ont du mal à s'approvisionner directement auprès de
producteurs locaux de confiance.

\paragraph{Gaspillage alimentaire dans la restauration.}
Le secteur de la restauration génère des surplus alimentaires considérables chaque jour.
En l'absence de canaux adaptés, ces invendus finissent généralement au rebut, alors qu'ils
pourraient être redistribués à prix réduit aux citoyens.

\paragraph{Absence de vérification et de traçabilité.}
Dans les transactions entre professionnels, il n'existe pas de mécanisme fiable de
vérification d'identité et de statut professionnel, ce qui génère une méfiance et freine
les échanges.

\section{Motivation du projet}

La motivation principale de VitaChain est de proposer une solution technologique holistique
qui s'attaque simultanément à ces quatre problèmes à travers une architecture modulaire,
extensible et sécurisée. Les motivations spécifiques sont :

\begin{itemize}
  \item \textbf{Impact social et environnemental :} réduction du gaspillage alimentaire,
        optimisation de l'usage de l'eau et des engrais, création de liens directs entre
        producteurs et acheteurs ;
  \item \textbf{Modernisation agricole :} adoption de l'agriculture de précision grâce à
        des capteurs IoT accessibles aux agriculteurs marocains ;
  \item \textbf{Inclusion économique :} permettre aux petits agriculteurs de valoriser
        leur production en accédant directement au marché de la restauration ;
  \item \textbf{Défi technologique :} concevoir et déployer en production une architecture
        full-stack moderne intégrant IoT, RLS, workers asynchrones et authentification
        multi-rôle.
\end{itemize}

\section{Étude de l'existant}

Plusieurs solutions numériques existent sur le marché international et national, mais
aucune ne couvre l'intégralité du périmètre de VitaChain.

\begin{table}[H]
  \centering
  \caption{Comparaison des solutions existantes}
  \label{tab:existant}
  \small
  \begin{tabularx}{\linewidth}{Xcccc}
    \toprule
    \textbf{Solution} & \textbf{IoT Agricole} & \textbf{Marketplace B2B}
                      & \textbf{Anti-gaspillage} & \textbf{Marché marocain} \\
    \midrule
    Agromonitoring    & \checkmark & $\times$ & $\times$ & $\times$ \\
    Too Good To Go    & $\times$   & $\times$ & \checkmark & Partiel  \\
    AgriApp           & $\times$   & \checkmark & $\times$ & $\times$ \\
    Fellah Trade      & $\times$   & \checkmark & $\times$ & \checkmark \\
    \rowcolor{rowHead}
    \textbf{VitaChain} & \checkmark & \checkmark & \checkmark & \checkmark \\
    \bottomrule
  \end{tabularx}
\end{table}

Les limites des solutions existantes sont les suivantes :

\begin{itemize}
  \item Les solutions IoT agricoles (Agromonitoring, Farmonaut) sont conçues pour des
        exploitations industrielles avec un coût d'entrée élevé.
  \item Les marketplaces B2B agricoles marocaines (Fellah Trade) n'intègrent pas de
        mécanisme de vérification KYC ni de suivi de commande en temps réel.
  \item Les applications anti-gaspillage (Too Good To Go) ne sont pas adaptées au
        contexte marocain.
  \item Aucune solution ne propose un tableau de bord IoT couplé à une place de marché
        et une fonctionnalité anti-gaspillage dans une même plateforme.
\end{itemize}

\section{Objectifs du projet}

Les objectifs de VitaChain sont organisés en quatre axes.

\paragraph{Axe 1 --- Module Katara (Surveillance IoT).}
\begin{itemize}
  \item Permettre aux agriculteurs de créer et gérer leurs parcelles agricoles
        géoréférencées.
  \item Intégrer des capteurs ESP32 envoyant des données de sol (humidité, température,
        pH, conductivité électrique) en temps réel.
  \item Afficher des graphiques d'historique et des alertes déclenchées sur des seuils
        configurables.
  \item Fournir des diagnostics agronomiques automatisés via intelligence artificielle
        (Google Gemini).
  \item Intégrer des données météorologiques (OpenWeatherMap) et satellitaires
        (Sentinel-2 NDVI).
\end{itemize}

\paragraph{Axe 2 --- Module FarMarket (Marketplace Agricole).}
\begin{itemize}
  \item Permettre aux agriculteurs vérifiés de publier des annonces de vente de produits
        agricoles avec photos.
  \item Permettre aux restaurateurs de parcourir le catalogue, filtrer par
        région/produit/prix, et passer des commandes.
  \item Gérer le cycle de vie complet d'une commande (acceptation, suivi, livraison,
        paiement).
  \item Protéger la confidentialité des restaurateurs vis-à-vis des agriculteurs
        (anonymisation).
  \item Intégrer un système de notation agriculteur par les restaurateurs.
\end{itemize}

\paragraph{Axe 3 --- Module SecondServe (Anti-gaspillage).}
\begin{itemize}
  \item Permettre aux restaurants de publier leurs surplus alimentaires à prix réduit.
  \item Permettre aux citoyens de réserver et commander ces offres.
  \item Assurer la continuité de session entre VitaChain et l'application SecondServe
        via SSO (magic-link OTP).
\end{itemize}

\paragraph{Axe 4 --- Gouvernance et Administration.}
\begin{itemize}
  \item Mettre en place un processus de vérification KYC pour les professionnels.
  \item Permettre aux administrateurs de gérer les utilisateurs, les rôles et les
        contenus de la plateforme.
\end{itemize}

\section{Périmètre du projet}

\paragraph{Inclus dans le périmètre :}
inscription multi-rôle ; vérification KYC ; module Katara complet ; module FarMarket
complet ; module SecondServe ; panneau d'administration ; infrastructure de production
(Docker, NGINX, TLS, sauvegardes, monitoring).

\paragraph{Hors périmètre :}
module BotaBa9a (architecture définie, implémentation différée) ; application mobile
native ; intégration de passerelles de paiement réelles ; gestion logistique propre.

\begin{figure}[H]
\centering
\begin{tikzpicture}
  % ── Boîtes des modules (empilées) ──
  \fill[blue!75, rounded corners=6pt]    (-5.4, 4.4) rectangle (5.4, 5.8);
  \fill[green!50!black, rounded corners=6pt] (-5.4, 2.6) rectangle (5.4, 4.0);
  \fill[orange!70, rounded corners=6pt]  (-5.4, 0.8) rectangle (5.4, 2.2);
  \fill[gray!60, rounded corners=6pt]    (-5.4,-1.0) rectangle (5.4, 0.4);

  % ── Labels modules ──
  \node[font=\bfseries\small, text=white] at (0, 5.1)
    {Katara --- Surveillance IoT agricole};
  \node[font=\bfseries\small, text=white] at (0, 3.3)
    {FarMarket --- Marketplace Agricole B2B};
  \node[font=\bfseries\small, text=white] at (0, 1.5)
    {SecondServe --- Anti-gaspillage alimentaire};
  \node[font=\bfseries\small, text=white] at (0,-0.3)
    {Administration};

  % ── Boîtes acteurs (droite) ──
  \fill[blue!15]  (6.0, 4.65) rectangle (9.0, 5.55);
  \draw[blue!60, rounded corners=3pt] (6.0, 4.65) rectangle (9.0, 5.55);
  \node[font=\footnotesize, text=blue!80] at (7.5, 5.1) {Agriculteur + ESP32};

  \fill[green!15]  (6.0, 2.85) rectangle (9.0, 3.75);
  \draw[green!45!black, rounded corners=3pt] (6.0, 2.85) rectangle (9.0, 3.75);
  \node[font=\footnotesize, text=green!45!black] at (7.5, 3.3) {Restaurateur};

  \fill[orange!15] (6.0, 1.05) rectangle (9.0, 1.95);
  \draw[orange!70, rounded corners=3pt] (6.0, 1.05) rectangle (9.0, 1.95);
  \node[font=\footnotesize, text=orange!80] at (7.5, 1.5) {Citoyen};

  \fill[gray!15]  (6.0,-0.75) rectangle (9.0, 0.15);
  \draw[gray!55, rounded corners=3pt] (6.0,-0.75) rectangle (9.0, 0.15);
  \node[font=\footnotesize, text=gray!70] at (7.5,-0.3) {Administrateur};

  % ── Flèches ──
  \draw[->, thick, blue!60]        (5.4, 5.1) -- (6.0, 5.1);
  \draw[->, thick, green!40!black] (5.4, 3.3) -- (6.0, 3.3);
  \draw[->, thick, orange!70]      (5.4, 1.5) -- (6.0, 1.5);
  \draw[->, thick, gray!55]        (5.4,-0.3) -- (6.0,-0.3);

  % ── Titre ──
  \node[font=\large\bfseries, text=blue!80] at (0, 6.6)
    {Plateforme VitaChain};
  \node[font=\small\itshape, text=gray!70] at (0, 6.1)
    {Architecture modulaire --- périmètre fonctionnel};
\end{tikzpicture}
\caption{Périmètre fonctionnel de VitaChain --- diagramme en couches des quatre modules}
\label{fig:1.1}
\end{figure}

\section{Acteurs}

VitaChain implique quatre types d'acteurs principaux et un acteur système.

\begin{figure}[H]
    \centering
    \includegraphics[width=0.8\textwidth]{fig1_2_acteurs}
    \caption{Diagramme des acteurs}
    \label{fig:fig1_2}
\end{figure}

\begin{table}[H]
  \centering
  \caption{Acteurs humains de VitaChain}
  \label{tab:acteurs}
  \begin{tabularx}{\linewidth}{llX}
    \toprule
    \textbf{Acteur} & \textbf{Rôle} & \textbf{Description} \\
    \midrule
    Agriculteur   & FARMER     & Producteur agricole. Crée des parcelles, associe des
                                 capteurs, publie des annonces, reçoit des commandes.
                                 Doit être vérifié (KYC). \\
    Restaurateur  & RESTAURANT & Acheteur professionnel. Consulte le catalogue, passe
                                 des commandes, évalue les agriculteurs, publie des
                                 offres surplus. Doit être vérifié. \\
    Citoyen       & CITIZEN    & Consommateur final. Consulte et commande des offres
                                 surplus via SecondServe. \\
    Administrateur& ADMIN      & Gestionnaire. Approuve les KYC, modère les contenus,
                                 gère les utilisateurs et les rôles. \\
    \midrule
    Dispositif ESP32 & ---     & Capteur IoT. Envoie des données télémétriques via HTTP
                                 (authentification par clé d'API). \\
    \bottomrule
  \end{tabularx}
\end{table}

\section{Besoins fonctionnels}

\subsection*{Module Authentification \& KYC}
\begin{enumerate}[label=BF-AUTH-\arabic*]
  \item L'utilisateur peut s'inscrire avec un email, un mot de passe et un rôle.
  \item L'utilisateur peut se connecter et obtenir un jeton JWT portant son rôle et
        son statut de vérification.
  \item Un professionnel peut soumettre des documents KYC (CIN, RC, carte agricole).
  \item L'administrateur peut approuver ou rejeter une demande KYC.
  \item Les actions professionnelles sont bloquées pour les profils non vérifiés.
\end{enumerate}

\subsection*{Module Katara}
\begin{enumerate}[label=BF-KAT-\arabic*]
  \item L'agriculteur peut créer, modifier et supprimer des parcelles géoréférencées
        (polygone GeoJSON).
  \item L'agriculteur peut associer un dispositif ESP32 à une parcelle et recevoir une
        clé d'API unique.
  \item Le dispositif ESP32 peut envoyer des données de télémétrie (humidité sol,
        température sol, pH, conductivité, batterie).
  \item L'agriculteur peut visualiser les données en temps réel et l'historique sous
        forme de graphiques (24h, 7j, 30j).
  \item L'agriculteur peut configurer des seuils d'alerte par métrique et recevoir un
        email en cas de dépassement.
  \item L'agriculteur peut demander un diagnostic agronomique IA et le consulter.
  \item L'agriculteur peut consulter les prévisions météorologiques pour sa parcelle.
  \item L'agriculteur peut consulter les images NDVI Sentinel-2 pour sa parcelle.
  \item Le système détecte automatiquement les dispositifs hors ligne
        (silence $>$ 1 heure).
\end{enumerate}

\subsection*{Module FarMarket}
\begin{enumerate}[label=BF-FAR-\arabic*]
  \item L'agriculteur vérifié peut créer une annonce avec titre, description, type de
        produit, prix, quantité, région et jusqu'à 5 photos.
  \item Le restaurateur vérifié peut parcourir le catalogue et filtrer par région,
        type de produit et fourchette de prix.
  \item Le restaurateur peut passer une commande multi-articles.
  \item L'agriculteur reçoit les commandes entrantes de manière anonymisée.
  \item L'agriculteur peut accepter/rejeter et mettre à jour le statut de chaque
        article.
  \item Le restaurateur peut confirmer un paiement en espèces à la livraison (COD).
  \item Le restaurateur peut laisser une évaluation (1--5 étoiles + commentaire) après
        livraison.
  \item Les annonces expirent automatiquement après 7 jours.
\end{enumerate}

\subsection*{Module SecondServe}
\begin{enumerate}[label=BF-SEC-\arabic*]
  \item Le restaurateur peut publier des offres de surplus alimentaires.
  \item Le citoyen peut consulter les offres, les filtrer par ville et passer une
        commande.
  \item L'utilisateur VitaChain peut accéder à SecondServe via un lien magique SSO.
\end{enumerate}

\subsection*{Module Administration}
\begin{enumerate}[label=BF-ADM-\arabic*]
  \item L'administrateur peut consulter et traiter la file d'attente des demandes KYC.
  \item L'administrateur peut lister, filtrer, bannir/débannir et changer le rôle
        d'un utilisateur.
  \item L'administrateur peut consulter le tableau de bord FarMarket.
\end{enumerate}

\section{Besoins non fonctionnels}

\begin{table}[H]
  \centering
  \caption{Besoins non fonctionnels de VitaChain}
  \label{tab:bnf}
  \begin{tabularx}{\linewidth}{lX}
    \toprule
    \textbf{Catégorie} & \textbf{Exigence} \\
    \midrule
    Performance
      & L'endpoint d'ingestion télémétrique doit répondre en moins de 50~ms (p50). \\
    Disponibilité
      & Disponibilité cible de 99,5~\% (24h/24, 7j/7). \\
    Sécurité
      & Toutes les communications chiffrées (HTTPS/TLS 1.3). Données protégées par
        RLS au niveau de la base de données. \\
    Authentification
      & Jetons JWT de durée de vie 1~heure avec rafraîchissement automatique. \\
    Isolation des données
      & Matrice RLS 22 cellules (rôle $\times$ table $\times$ verbe). \\
    Scalabilité
      & Architecture Docker/Gunicorn/NGINX permettant une montée en charge
        horizontale. \\
    Maintenabilité
      & Tests pytest, pipelines CI/CD GitHub Actions (5 jobs). \\
    Sauvegarde
      & Sauvegarde nocturne pg\_dump + rclone vers Backblaze B2. \\
    Observabilité
      & Sentry (erreurs), Uptime Kuma (disponibilité). \\
    Internationalisation
      & Interface SecondServe en français et arabe. \\
    \bottomrule
  \end{tabularx}
\end{table}

\section{Contraintes techniques}

\paragraph{Contraintes d'infrastructure.}
Déploiement sur un VPS Linux avec ressources limitées. Utilisation de Supabase comme
backend-as-a-service. Conteneurisation obligatoire via Docker.

\paragraph{Contraintes de sécurité.}
La clé de service Supabase ne doit jamais être exposée côté client (AUTH-05). Toutes
les mutations de données doivent passer par les politiques RLS de PostgreSQL. Les clés
d'API des dispositifs IoT doivent être hachées avec bcrypt avant stockage. Aucune donnée
personnelle du restaurateur ne doit être visible par l'agriculteur.

\paragraph{Contraintes de développement.}
Le backend FastAPI doit respecter la structure de routeurs modulaires par domaine métier.
Le frontend Next.js doit utiliser les composants serveur (RSC) pour les lectures d'état
d'authentification. Les migrations de base de données doivent être idempotentes et
versionnées séquentiellement.

\paragraph{Contraintes matérielles (IoT).}
Microcontrôleur ESP32 avec connectivité Wi-Fi. Protocole HTTP/1.1 (POST JSON).
Authentification par en-têtes HTTP (\code{X-Device-Id} + \code{X-Device-Api-Key}).

\section{Planning prévisionnel}

Le développement de VitaChain a été planifié selon une méthodologie agile Scrum,
organisée en sprints de deux semaines.

\begin{figure}[H]
\centering
\begin{ganttchart}[
  hgrid,
  vgrid        = {*{1}{draw=gray!20, dashed}},
  x unit       = 0.56cm,
  y unit title = 0.72cm,
  y unit chart = 0.60cm,
  title label font = \tiny\bfseries,
  bar label font   = \footnotesize,
  title height = 1,
  bar height   = 0.60,
]{1}{22}
  \gantttitle{Planning prévisionnel VitaChain --- 22 semaines}{22} \\
  \gantttitlelist{1,...,22}{1} \\
  \ganttbar[bar/.append style={fill=gray!55}]
    {Sprint 0 --- Infrastructure}{1}{2} \\
  \ganttbar[bar/.append style={fill=purple!55}]
    {Sprint 1 --- Authentification \& KYC}{3}{4} \\
  \ganttbar[bar/.append style={fill=blue!80}]
    {Sprint 2 --- Katara : Parcelles}{5}{6} \\
  \ganttbar[bar/.append style={fill=blue!65}]
    {Sprint 3 --- Katara : Tableaux de bord}{7}{8} \\
  \ganttbar[bar/.append style={fill=blue!45}]
    {Sprint 4 --- Katara : Alertes \& IA}{9}{10} \\
  \ganttbar[bar/.append style={fill=green!55!black}]
    {Sprint 5 --- FarMarket : Annonces}{11}{12} \\
  \ganttbar[bar/.append style={fill=green!45!black}]
    {Sprint 6 --- FarMarket : Commandes}{13}{14} \\
  \ganttbar[bar/.append style={fill=green!35!black}]
    {Sprint 7 --- FarMarket : Paiement \& Avis}{15}{16} \\
  \ganttbar[bar/.append style={fill=orange!70}]
    {Sprint 8 --- SecondServe}{17}{18} \\
  \ganttbar[bar/.append style={fill=red!50}]
    {Sprint 9 --- Administration}{19}{20} \\
  \ganttbar[bar/.append style={fill=gray!40}]
    {Sprint 10 --- Tests \& Livraison}{21}{22}
\end{ganttchart}
\caption{Diagramme de Gantt prévisionnel --- planning des sprints VitaChain (22 semaines)}
\label{fig:1.3}
\end{figure}

\begin{table}[H]
  \centering
  \caption{Planning prévisionnel des sprints}
  \label{tab:planning}
  \begin{tabularx}{\linewidth}{llcX}
    \toprule
    \textbf{Phase} & \textbf{Sprints} & \textbf{Durée} & \textbf{Contenu} \\
    \midrule
    Infrastructure & Sprint 0   & 2 sem. & VPS, Supabase, CI/CD, Docker, NGINX, HTTPS \\
    Authentification & Sprint 1 & 2 sem. & Inscription, connexion, rôles, JWT, RLS, KYC \\
    Katara IoT  & Sprints 2--4  & 6 sem. & Parcelles, dispositifs, télémétrie, alertes, IA \\
    FarMarket   & Sprints 5--7  & 6 sem. & Annonces, catalogue, commandes, paiements,
                                           évaluations \\
    SecondServe & Sprint 8      & 2 sem. & Offres surplus, commandes citoyens, SSO \\
    Administration & Sprint 9   & 2 sem. & File KYC, gestion utilisateurs, tableau de bord \\
    Tests \& Livraison & Sprint 10 & 2 sem. & Tests de charge, corrections, documentation \\
    \bottomrule
  \end{tabularx}
\end{table}

\section{Conclusion}

Ce premier chapitre a permis d'établir le cadre complet du projet VitaChain : son contexte
dans la transformation numérique agricole marocaine, les problèmes structurels qu'il
adresse, ses objectifs fonctionnels et non fonctionnels, ainsi que les contraintes
techniques qui guident sa conception. La plateforme couvre quatre modules complémentaires
--- Katara, FarMarket, SecondServe et Administration --- répondant chacun à un besoin
identifié dans la chaîne agroalimentaire. Le chapitre suivant détaille l'analyse et la
conception de ces modules.

% =============================================================
% CHAPITRE 2 — ANALYSE ET CONCEPTION
% =============================================================
\chapter{Analyse et Conception}

\section{Introduction}

Ce chapitre présente la démarche d'analyse et de conception adoptée pour VitaChain.
Nous commençons par identifier les sous-modules du système et les cas d'utilisation
associés à chaque acteur, puis nous décrivons l'architecture logicielle et le modèle
de données. La conception s'appuie sur des patterns architecturaux éprouvés pour
garantir la sécurité, la modularité et la maintenabilité de la plateforme.

\section{Analyse}

\subsection{Identification des sous-modules}

L'analyse fonctionnelle de VitaChain révèle une architecture modulaire organisée autour
de cinq sous-systèmes indépendants mais interconnectés.

\begin{figure}[H]
\centering
\resizebox{0.97\linewidth}{!}{%
\begin{tikzpicture}[
  racine/.style={draw=blue!70, fill=blue!80, text=white, rounded corners=5pt,
                 font=\bfseries, minimum width=4.2cm, minimum height=1.0cm, align=center},
  module/.style={rounded corners=4pt, font=\small\bfseries,
                 minimum width=3.0cm, minimum height=0.9cm, align=center},
  comp/.style={draw=gray!50, fill=gray!8, rounded corners=3pt, font=\footnotesize,
               minimum width=2.8cm, minimum height=0.7cm, align=center},
  arr/.style={->, thick, gray!55, shorten >=2pt, shorten <=2pt}
]
  \node[racine] (vc) at (0,0) {VitaChain};

  \node[module, draw=purple!60, fill=purple!10]         (auth) at (-10.5,-2.5) {Auth \& KYC};
  \node[module, draw=blue!60,   fill=blue!10]           (kat)  at (-5.5, -2.5) {Katara (IoT)};
  \node[module, draw=green!60!black, fill=green!10]     (far)  at  (0.0, -2.5) {FarMarket};
  \node[module, draw=orange!70, fill=orange!10]         (ss)   at  (5.5, -2.5) {SecondServe};
  \node[module, draw=gray!60,   fill=gray!12]           (adm)  at (10.5,-2.5) {Administration};

  \foreach \m in {auth,kat,far,ss,adm}
    \draw[arr] (vc.south) -- ++(0,-0.55) -| (\m.north);

  \node[comp] (c11) at (-10.5,-4.3) {Inscription};
  \node[comp] (c12) at (-10.5,-5.2) {Connexion / JWT};
  \node[comp] (c13) at (-10.5,-6.1) {Soumission KYC};

  \node[comp] (c21) at (-5.5,-4.3) {Parcelles};
  \node[comp] (c22) at (-5.5,-5.2) {Dispositifs ESP32};
  \node[comp] (c23) at (-5.5,-6.1) {Télémétrie};
  \node[comp] (c24) at (-5.5,-7.0) {Alertes / Seuils};
  \node[comp] (c25) at (-5.5,-7.9) {Diagnostics IA};

  \node[comp] (c31) at (0,-4.3) {Annonces};
  \node[comp] (c32) at (0,-5.2) {Catalogue};
  \node[comp] (c33) at (0,-6.1) {Commandes};
  \node[comp] (c34) at (0,-7.0) {Évaluations};

  \node[comp] (c41) at (5.5,-4.3) {Offres surplus};
  \node[comp] (c42) at (5.5,-5.2) {Commandes COD};
  \node[comp] (c43) at (5.5,-6.1) {SSO Magic Link};

  \node[comp] (c51) at (10.5,-4.3) {File KYC};
  \node[comp] (c52) at (10.5,-5.2) {Utilisateurs};
  \node[comp] (c53) at (10.5,-6.1) {Modération};

  \foreach \c in {c11,c12,c13}
    \draw[arr] (auth.south) -- ++(0,-0.3) -| (\c.north);
  \foreach \c in {c21,c22,c23,c24,c25}
    \draw[arr] (kat.south)  -- ++(0,-0.3) -| (\c.north);
  \foreach \c in {c31,c32,c33,c34}
    \draw[arr] (far.south)  -- ++(0,-0.3) -| (\c.north);
  \foreach \c in {c41,c42,c43}
    \draw[arr] (ss.south)   -- ++(0,-0.3) -| (\c.north);
  \foreach \c in {c51,c52,c53}
    \draw[arr] (adm.south)  -- ++(0,-0.3) -| (\c.north);
\end{tikzpicture}}
\caption{Diagramme de décomposition modulaire de VitaChain}
\label{fig:2.1}
\end{figure}

\paragraph{Sous-module 1 --- Auth \& KYC.}
Gère l'inscription, la connexion, la gestion des sessions JWT et le processus de
vérification d'identité professionnelle (KYC). Ce sous-module est transversal : il
conditionne l'accès aux fonctionnalités de tous les autres modules.

Composants : service d'inscription multi-rôle ; générateur de jetons JWT avec claims
personnalisés (\code{user\_role}, \code{verification\_status}) ; gestionnaire de
documents KYC (upload sécurisé, soumission, révision admin) ; Hook Supabase Custom
Access Token pour la propagation des rôles dans les JWT.

\paragraph{Sous-module 2 --- Katara (IoT Agricole).}
Gère le cycle de vie des parcelles et des dispositifs IoT, l'ingestion et la
visualisation des données télémétriques, ainsi que les services à valeur ajoutée
(alertes, diagnostics IA, météo, NDVI).

Composants : registre de parcelles (CRUD géoréférencé) ; gestionnaire de dispositifs
ESP32 (appairage, clé d'API, historique) ; pipeline d'ingestion télémétrique
(hot-path $<$ 50~ms) ; service de visualisation (graphiques temps réel et historique) ;
moteur d'alertes sur seuils (LISTEN/NOTIFY + worker email) ; orchestrateur de
diagnostics IA (Gemini + OWM + Sentinel-2) ; détecteur de dispositifs hors ligne
(CRON worker).

\paragraph{Sous-module 3 --- FarMarket (Marketplace).}
Gère la publication et la consultation d'annonces agricoles, le cycle de vie des
commandes entre restaurateurs et agriculteurs, et le système de notation.

Composants : gestionnaire d'annonces (CRUD, photos, expiration) ; catalogue public
(filtres, pagination, recherche textuelle pg\_trgm) ; moteur de commandes
(multi-articles, anonymisation, suivi de statut) ; service de paiement (COD
simulation, audit log) ; système d'évaluation (1--5~étoiles, contrainte achat
vérifié).

\paragraph{Sous-module 4 --- SecondServe (Anti-gaspillage).}
Application distincte (Vite/React) connectée au même backend Supabase, gérant les
offres surplus des restaurants et les commandes des citoyens.

Composants : gestionnaire d'offres surplus (publication restaurant, visibilité
conditionnelle) ; catalogue citoyen (filtres par ville) ; moteur de commandes citoyen
(checkout, suivi) ; SSO inter-application (magic-link OTP depuis VitaChain).

\paragraph{Sous-module 5 --- Administration.}
File d'attente KYC (revue, approbation, rejet) ; gestion des utilisateurs (liste,
rôle, ban) ; tableau de bord FarMarket (statistiques, modération).

\subsection{Sous Modules}

\subsubsection{Authentification et verification KYC}
\begin{figure}[H]
    \centering
    \includegraphics[width=0.8\textwidth]{fig2_2_uc_auth.jpeg}
    \caption{Diagramme d'authentification et verification KYC}
    \label{fig:fig1_2}
\end{figure}
\paragraph{UC-AUTH-01 : Inscription d'un professionnel.}
\textit{Acteur :} Agriculteur / Restaurateur. \textit{Précondition :} l'utilisateur
n'a pas de compte VitaChain.
\begin{enumerate}
  \item L'utilisateur accède à la page d'inscription.
  \item Il saisit son email, son mot de passe et sélectionne son rôle.
  \item Le système crée un compte Supabase Auth et un profil \code{profiles} avec
        statut PENDING.
  \item L'utilisateur est redirigé vers son tableau de bord avec accès limité.
\end{enumerate}
\textit{Postcondition :} compte créé avec \code{verification\_status = PENDING}.

\paragraph{UC-AUTH-02 : Soumission KYC.}
\textit{Acteur :} Agriculteur / Restaurateur (PENDING).
\begin{enumerate}
  \item L'utilisateur accède à la section \og Vérification \fg{}.
  \item Il télécharge ses documents (CIN, RC ou carte agricole) dans Supabase Storage.
  \item Il soumet sa demande via \api{POST /api/v1/kyc/submit}.
  \item L'administrateur reçoit une notification et la demande apparaît dans la file
        d'attente KYC.
\end{enumerate}
\textit{Postcondition :} demande KYC créée avec statut EN\_ATTENTE.

\paragraph{UC-AUTH-03 : Approbation KYC par l'administrateur.}
\textit{Acteur :} Administrateur.
\begin{enumerate}
  \item L'administrateur consulte la file d'attente.
  \item Il examine les documents soumis.
  \item Il approuve ou rejette la demande.
  \item Le système met à jour \code{profiles.verification\_status} à VERIFIED ou
        REJECTED.
  \item Un email de notification est envoyé via Brevo.
\end{enumerate}
\textit{Postcondition :} l'utilisateur peut réaliser des actions professionnelles.

\begin{figure}[H]
    \centering
    \includegraphics[width=0.8\textwidth]{gestion des utilisateur.jpeg}
    \caption{Diagrammes de gestion des utilisateurs}
    \label{fig:fig_gestion_utilisateurs}
\end{figure}

\subsubsection{Gestion des parcelles}
\begin{figure}[H]
    \centering
    \includegraphics[width=0.8\textwidth]{fig2_3a_uc_katara_parcelles.jpeg}
    \caption{Diagramme des parcelles katara}
    \label{fig:fig1_2}
\end{figure}
\subsubsection{Gestion des dispositives IOT }
\begin{figure}[H]
    \centering
    \includegraphics[width=0.8\textwidth]{fig2_3b_uc_katara_iot.jpeg}
    \caption{Diagramme iot katara}
    \label{fig:fig1_2}
\end{figure}
\subsubsection{Gestion des telemetries}
\begin{figure}[H]
    \centering
    \includegraphics[width=0.8\textwidth]{fig2_3c_uc_katara_telemetrie.jpeg}
    \caption{Diagramme de telemetrie katara}
    \label{fig:fig1_2}
\end{figure}
\subsubsection{Gestion du diagnostique ai}
\begin{figure}[H]
    \centering
    \includegraphics[width=0.8\textwidth]{fig2_3d_uc_katara_diagnostic.jpeg}
    \caption{Diagramme de diagnostique katara}
    \label{fig:fig1_2}
\end{figure}

\paragraph{UC-KAT-01 : Enregistrement d'une parcelle.}
\textit{Acteur :} Agriculteur (VERIFIED).
\begin{enumerate}
  \item L'agriculteur clique sur \og Nouvelle parcelle \fg{} et renseigne le nom,
        le type de culture et la surface.
  \item Il dessine le contour sur la carte Leaflet (polygone GeoJSON).
  \item Il soumet le formulaire via \api{POST /api/v1/katara/parcels}.
\end{enumerate}
\textit{Postcondition :} parcelle créée avec géométrie GeoJSON stockée.

\paragraph{UC-KAT-02 : Appairage d'un dispositif ESP32.}
\textit{Acteur :} Agriculteur (VERIFIED).
\begin{enumerate}
  \item L'agriculteur sélectionne une parcelle et clique sur \og Associer un capteur \fg{}.
  \item Il saisit l'identifiant du dispositif (format \code{ESP-KAT-NNN}).
  \item Le backend génère une clé d'API unique (\code{vk\_<32 hex>}), la hache
        (bcrypt) et la stocke.
  \item La clé en clair est affichée \textbf{une seule fois} à l'agriculteur.
\end{enumerate}
\textit{Postcondition :} dispositif associé avec statut PENDING $\to$ ACTIVE à la
première télémétrie.

\paragraph{UC-KAT-03 : Ingestion de données télémétriques.}
\textit{Acteur :} Dispositif ESP32.
\begin{enumerate}
  \item L'ESP32 envoie une requête HTTP POST avec \code{X-Device-Id} et
        \code{X-Device-Api-Key}.
  \item Le backend appelle la fonction SQL \code{m1\_katara\_ingest()} (SECURITY
        DEFINER).
  \item La fonction vérifie la clé (bcrypt), insère la mesure, met à jour
        \code{last\_seen} et envoie un NOTIFY.
  \item Le worker de seuils écoute le NOTIFY et vérifie les dépassements.
  \item Si un seuil est franchi, un email d'alerte est envoyé à l'agriculteur.
\end{enumerate}
\textit{Postcondition :} mesure enregistrée. Durée totale $<$ 50~ms (p50).

\paragraph{UC-KAT-04 : Demande de diagnostic IA.}
\textit{Acteur :} Agriculteur (VERIFIED).
\begin{enumerate}
  \item L'agriculteur clique sur \og Demander un diagnostic \fg{}.
  \item Le backend crée un enregistrement de diagnostic avec statut PENDING.
  \item Un worker asynchrone collecte les données : dernières mesures, prévisions
        météo (OWM), images NDVI (Sentinel-2).
  \item Les données sont envoyées à l'API Gemini.
  \item Le résumé diagnostique est stocké ; statut passe à COMPLETED.
  \item Un email de notification est envoyé.
\end{enumerate}
\textit{Postcondition :} diagnostic consultable avec résumé en langage naturel.

\subsubsection{Gestion des commandes}

\begin{figure}[H]
    \centering
    \includegraphics[width=0.8\textwidth]{fig2_4b_uc_farmarket_commandes.jpeg}
    \caption{Diagramme des commandes Farmarket}
    \label{fig:fig1_2}
\end{figure}
\subsubsection{Gestion des evaluation}
\begin{figure}[H]
    \centering
    \includegraphics[width=0.8\textwidth]{fig2_4c_uc_farmarket_eval.jpeg}
    \caption{Diagramme de evaluation Farmarket}
    \label{fig:fig1_2}
\end{figure}

\paragraph{UC-FAR-01 : Publication d'une annonce agricole.}
\textit{Acteur :} Agriculteur (VERIFIED).
\begin{enumerate}
  \item L'agriculteur renseigne le titre, la description, le type de produit, le prix
        (MAD/kg), la quantité et la région.
  \item Il télécharge jusqu'à 5 photos du produit.
  \item Il soumet l'annonce via \api{POST /api/v1/farmarket/ads}.
\end{enumerate}
\textit{Postcondition :} annonce ACTIVE visible dans le catalogue, expiration à J+7.

\paragraph{UC-FAR-02 : Passage d'une commande.}
\textit{Acteur :} Restaurateur (VERIFIED).
\begin{enumerate}
  \item Le restaurateur ajoute des articles au panier depuis plusieurs agriculteurs.
  \item Il sélectionne le mode de paiement (COD ou PSP\_TRANSFER) et la région de
        livraison.
  \item Le système calcule les frais logistiques selon la région.
  \item La commande est créée avec instantanés de prix (\code{unit\_price\_snapshot})
        pour chaque article.
  \item Chaque agriculteur reçoit une notification de nouvelle commande.
\end{enumerate}
\textit{Postcondition :} commande PENDING ; les agriculteurs voient leurs articles
entrants anonymisés.

\paragraph{UC-FAR-03 : Suivi et livraison d'une commande.}
\textit{Acteurs :} Agriculteur (par article), Restaurateur (vue globale).
\begin{enumerate}
  \item L'agriculteur consulte ses articles entrants (vue anonymisée).
  \item Il accepte ou rejette chaque article.
  \item Le statut de l'article évolue : PENDING $\to$ ACCEPTED $\to$ PICKED\_UP
        $\to$ IN\_TRANSIT $\to$ DELIVERED.
  \item Le restaurateur confirme le paiement COD après réception.
\end{enumerate}
\textit{Postcondition :} commande DELIVERED, paiement confirmé.

\subsubsection{Gestion des offres}

\begin{figure}[H]
    \centering
    \includegraphics[width=0.8\textwidth]{fig2_5a_uc_ss_offres.jpeg}
    \caption{Diagramme des offres SecondServe}
    \label{fig:fig1_2}
\end{figure}

\paragraph{UC-SEC-01 : Accès SSO depuis VitaChain.}
\textit{Acteur :} Restaurateur / Citoyen.
\begin{enumerate}
  \item L'utilisateur clique sur \og Accéder à SecondServe \fg{} depuis VitaChain.
  \item Le backend génère un OTP Supabase à usage unique via
        \api{POST /api/v1/secondserve/handoff}.
  \item L'utilisateur est redirigé vers SecondServe avec le token dans l'URL.
  \item SecondServe échange le token contre une session Supabase valide.
\end{enumerate}
\textit{Postcondition :} utilisateur connecté à SecondServe sans ressaisie
d'identifiants.

\subsection{Diagramme de classes d'analyse}

L'analyse orientée objet du domaine de VitaChain identifie les classes principales
suivantes.

\begin{figure}[H]
    \centering
    \includegraphics[width=0.8\textwidth]{fig2_6_classes.png}
    \caption{Diagramme de classes d'analyse}
    \label{fig:fig2_6_classes}
\end{figure}

\paragraph{Classes du domaine Auth.}
\code{User} (email, password\_hash, created\_at) --- géré par Supabase Auth ;
\code{Profile} (id, role, verification\_status, full\_name, phone, locale) ;
\code{KycDocument} (id, user\_id, type, storage\_path, status, reviewed\_at,
reviewed\_by).

\paragraph{Classes du domaine Katara.}
\code{Parcel} (id, farmer\_id, name, geojson, crop\_type, surface\_area\_ha) ;
\code{Device} (id, device\_id, parcel\_id, api\_key\_hash, status, last\_seen) ;
\code{TelemetryReading} (id, device\_id, parcel\_id, soil\_moisture, soil\_temperature,
soil\_ph, soil\_conductivity, battery\_level, recorded\_at) ;
\code{AlertThreshold} (id, parcel\_id, metric, min\_value, max\_value, enabled) ;
\code{AiDiagnostic} (id, parcel\_id, status, diagnostic\_summary).

\paragraph{Classes du domaine FarMarket.}
\code{FarmarketAd} (id, farmer\_id, title, product\_type, price\_mad, quantity\_kg,
region, photo\_paths, status, expires\_at) ;
\code{Order} (id, restaurant\_id, status, delivery\_region, subtotal\_mad,
logistics\_fee\_mad, total\_mad, payment\_method, payment\_status) ;
\code{OrderItem} (id, order\_id, ad\_id, farmer\_id, quantity\_kg, unit\_price\_mad,
line\_total\_mad, status) ;
\code{FarmerRating} (id, farmer\_id, restaurant\_id, rating, review).

\paragraph{Relations clés.}
\code{Profile} 1---N \code{Parcel} ;
\code{Parcel} 1---N \code{Device} (1 actif max) ;
\code{Device} 1---N \code{TelemetryReading} ;
\code{Profile} (RESTAURANT) 1---N \code{Order} ;
\code{Order} 1---N \code{OrderItem} ;
\code{OrderItem} N---1 \code{FarmarketAd}.

\section{Conception}

\subsection{Génération du modèle physique de données (MPD)}

Le modèle physique de données est implémenté en PostgreSQL~17 via Supabase. Il comprend
12 tables principales, 5 vues, 8 types ENUM et plusieurs fonctions et triggers de
sécurité.

\begin{figure}[H]
    \centering
    \includegraphics[width=\textwidth]{mdp.png}
    \caption{Modèle Physique de Données (MPD) complet}
    \label{fig:fig1_2}
\end{figure}

\begin{table}[H]
  \centering
  \caption{Types énumérés PostgreSQL de VitaChain}
  \label{tab:enums}
  \small
  \begin{tabularx}{\linewidth}{lX}
    \toprule
    \textbf{Type} & \textbf{Valeurs} \\
    \midrule
    \code{user\_role}
      & FARMER, RESTAURANT, CITIZEN, ADMIN \\
    \code{verification\_status}
      & PENDING, VERIFIED, REJECTED \\
    \code{device\_status}
      & PENDING, ACTIVE, OFFLINE, UNLINKED \\
    \code{m2\_farmarket\_ad\_status}
      & ACTIVE, EXPIRED, DELETED \\
    \code{m2\_farmarket\_order\_status}
      & PENDING, PARTIALLY\_ACCEPTED, ACCEPTED, REJECTED,
        IN\_PROGRESS, DELIVERED, CANCELLED \\
    \code{m2\_farmarket\_item\_status}
      & PENDING, ACCEPTED, REJECTED, PICKED\_UP, IN\_TRANSIT, DELIVERED \\
    \code{m2\_farmarket\_region}
      & 12 régions administratives du Maroc (Tanger-Tétouan-Al Hoceïma,
        Oriental, Fès-Meknès, Rabat-Salé-Kénitra, Béni Mellal-Khénifra,
        Casablanca-Settat, Marrakech-Safi, Drâa-Tafilalet, Souss-Massa,
        Guelmim-Oued Noun, Laâyoune-Sakia El Hamra,
        Dakhla-Oued Ed-Dahab) \\
    \bottomrule
  \end{tabularx}
\end{table}

\begin{table}[H]
  \centering
  \caption{Vues PostgreSQL de VitaChain}
  \label{tab:vues}
  \small
  \begin{tabularx}{\linewidth}{lX}
    \toprule
    \textbf{Vue} & \textbf{Description} \\
    \midrule
    \code{m1\_katara\_telemetry\_latest}
      & Dernière mesure par parcelle (DISTINCT ON). \\
    \code{m1\_katara\_parcel\_device\_history}
      & Historique de tous les dispositifs associés à une parcelle. \\
    \code{v\_farmer\_incoming\_items}
      & Articles entrants pour les agriculteurs avec
        \code{resto\_handle = sha256(restaurant\_id || farmer\_id)}. \\
    \code{v\_farmarket\_farmer\_public}
      & Profil public d'un agriculteur pour le catalogue. \\
    \code{v\_farmarket\_farmer\_rating\_stats}
      & Moyenne et nombre d'évaluations par agriculteur. \\
    \bottomrule
  \end{tabularx}
\end{table}

\subsection{Génération du script SQL}

Les migrations de base de données sont organisées en 50 fichiers SQL numérotés
séquentiellement dans le répertoire \code{db/migrations/}. Chaque fichier est idempotent
et suit la convention de nommage \code{NNNN\_description.sql}.

\begin{table}[H]
  \centering
  \caption{Extrait des migrations SQL de VitaChain}
  \label{tab:migrations}
  \small
  \begin{tabularx}{\linewidth}{clX}
    \toprule
    \textbf{N°} & \textbf{Fichier} & \textbf{Contenu} \\
    \midrule
    0001 & \code{0001\_extensions\_enums.sql}
         & Extensions pg\_trgm, pgcrypto ; types ENUM. \\
    0002 & \code{0002\_profiles.sql}
         & Table \code{profiles}, contraintes, index. \\
    0003 & \code{0003\_signup\_trigger.sql}
         & Trigger \code{handle\_new\_user()} pour création automatique
           du profil. \\
    0006 & \code{0006\_jwt\_role\_hook.sql}
         & Hook Supabase Custom Access Token (role dans JWT). \\
    0008 & \code{0008\_rls\_has\_role\_helper.sql}
         & Fonctions \code{has\_role()}, \code{is\_admin()} SECURITY
           DEFINER. \\
    0016 & \code{0016\_katara\_parcels.sql}
         & Table \code{m1\_katara\_parcels}, RLS, index. \\
    0017 & \code{0017\_katara\_devices.sql}
         & Table \code{m1\_katara\_devices}, contraintes unicité,
           \code{verify\_device\_api\_key()}. \\
    0018 & \code{0018\_katara\_telemetry.sql}
         & Table \code{m1\_katara\_telemetry}, fonction
           \code{m1\_katara\_ingest()}, FORCE RLS. \\
    0031 & \code{0031\_farmarket\_ads.sql}
         & Table \code{m2\_farmarket\_ads}, index trgm, storage policies. \\
    0040 & \code{0040\_farmarket\_orders.sql}
         & Tables commandes et articles, vue anonymisation. \\
    0046 & \code{0046\_farmarket\_ratings.sql}
         & Table évaluations, contrainte achat vérifié. \\
    0048 & \code{0048\_ss\_signup.sql}
         & Trigger SecondServe : bypass profil VitaChain pour
           \code{ss\_app} metadata. \\
    \bottomrule
  \end{tabularx}
\end{table}

Le listing~\ref{lst:ingest} présente la fonction SQL atomique d'ingestion
télémétrique, cœur du module Katara.

\begin{lstlisting}[style=sql, caption={Fonction \texttt{m1\_katara\_ingest()} --- ingestion atomique télémétrique},
                   label=lst:ingest]
CREATE OR REPLACE FUNCTION public.m1_katara_ingest(
  p_device_id         TEXT,
  p_api_key           TEXT,
  p_soil_moisture     NUMERIC,
  p_soil_temperature  NUMERIC,
  p_soil_ph           NUMERIC,
  p_soil_conductivity NUMERIC,
  p_battery_level     NUMERIC,
  p_recorded_at       TIMESTAMPTZ DEFAULT now()
) RETURNS void
LANGUAGE plpgsql SECURITY DEFINER AS $$
DECLARE
  v_parcel_id UUID;
  v_farmer_id UUID;
BEGIN
  -- Verification bcrypt de la cle d'API (temps constant)
  IF NOT public.verify_device_api_key(p_device_id, p_api_key) THEN
    RAISE EXCEPTION 'invalid_device_credentials';
  END IF;

  -- Recuperation du contexte
  SELECT parcel_id, farmer_id INTO v_parcel_id, v_farmer_id
    FROM m1_katara_devices
   WHERE device_id = p_device_id AND status = 'ACTIVE';

  -- Insertion idempotente
  INSERT INTO m1_katara_telemetry (
    device_id, parcel_id, farmer_id,
    soil_moisture, soil_temperature, soil_ph,
    soil_conductivity, battery_level, recorded_at)
  VALUES (
    p_device_id, v_parcel_id, v_farmer_id,
    p_soil_moisture, p_soil_temperature, p_soil_ph,
    p_soil_conductivity, p_battery_level, p_recorded_at)
  ON CONFLICT (device_id, recorded_at) DO NOTHING;

  -- Mise a jour last_seen + activation si PENDING
  UPDATE m1_katara_devices
     SET last_seen = now(),
         status = CASE WHEN status = 'PENDING'
                       THEN 'ACTIVE' ELSE status END
   WHERE device_id = p_device_id;

  -- Notification aux workers (threshold + diagnostic)
  PERFORM pg_notify('katara_telemetry',
    json_build_object(
      'device_id', p_device_id,
      'parcel_id', v_parcel_id,
      'farmer_id', v_farmer_id)::text);
END;
$$;
\end{lstlisting}

\subsection{Design Patterns utilisés}

\begin{table}[H]
  \centering
  \caption{Design Patterns utilisés dans VitaChain}
  \label{tab:patterns}
  \small
  \begin{tabularx}{\linewidth}{llX}
    \toprule
    \textbf{Pattern} & \textbf{Module} & \textbf{Description} \\
    \midrule
    Repository Pattern
      & Backend
      & \code{db.py} expose \code{service\_client()} et
        \code{user\_scoped\_client()} --- abstraction de l'accès Supabase. \\
    Router Pattern
      & FastAPI
      & Chaque domaine est encapsulé dans un \code{APIRouter} monté
        de manière modulaire. \\
    Dependency Injection
      & FastAPI
      & \code{get\_current\_user()}, \code{require\_role()},
        \code{require\_verified()} sont des dépendances FastAPI. \\
    Outbox Pattern
      & Notifications
      & Table \code{notifications\_outbox} : les mutations écrivent dans
        l'outbox, un worker lit et envoie les emails. \\
    CQRS (léger)
      & FarMarket
      & Séparation lecture (catalogue public, vue anonymisée) / écriture
        (commandes) avec politiques RLS distinctes. \\
    SECURITY DEFINER
      & Base de données
      & Fonctions critiques (\code{m1\_katara\_ingest},
        \code{verify\_device\_api\_key}, \code{has\_role}) s'exécutent
        avec les droits du propriétaire. \\
    Snapshot
      & FarMarket
      & Prix et identifiant agriculteur copiés au moment de la commande
        dans \code{order\_items}. \\
    Magic Link SSO
      & SecondServe
      & OTP Supabase à usage unique pour le handoff inter-application. \\
    Observer (LISTEN/NOTIFY)
      & Katara
      & La base publie des événements ; les workers Python s'abonnent
        via LISTEN. \\
    FORCE RLS
      & Télémétrie
      & \code{ALTER TABLE ... FORCE ROW LEVEL SECURITY} --- même le
        propriétaire de la table est soumis aux politiques. \\
    \bottomrule
  \end{tabularx}
\end{table}

\subsection{Architecture du module}

VitaChain suit une architecture \textbf{N-tiers} organisée en couches distinctes.

\begin{figure}[H]
\centering
\begin{tikzpicture}
  % ── Tier 1 : Présentation ──
  \fill[blue!10, rounded corners=6pt]  (-7.2,-0.7) rectangle (7.2, 0.7);
  \draw[blue!60, rounded corners=6pt]  (-7.2,-0.7) rectangle (7.2, 0.7);
  \node[font=\small\bfseries, text=blue!80] at (-6.0,0) {Tier 1};
  \node[font=\footnotesize, draw=blue!50,   fill=white, rounded corners=3pt,
        minimum width=3.0cm, minimum height=0.65cm] at (-2.5, 0) {Next.js 15 (RSC/SSR)};
  \node[font=\footnotesize, draw=orange!60, fill=orange!6, rounded corners=3pt,
        minimum width=3.0cm, minimum height=0.65cm] at (1.2,  0) {SecondServe (Vite 6)};
  \node[font=\footnotesize, draw=gray!50,   fill=gray!6,   rounded corners=3pt,
        minimum width=2.6cm, minimum height=0.65cm] at (4.7,  0) {NGINX 1.27};
  \node[font=\tiny\itshape, text=gray!70] at (0, 0.95) {Couche Présentation};

  % ── Tier 2 : API ──
  \fill[green!10, rounded corners=6pt]      (-7.2,-2.5) rectangle (7.2,-1.1);
  \draw[green!50!black, rounded corners=6pt] (-7.2,-2.5) rectangle (7.2,-1.1);
  \node[font=\small\bfseries, text=green!50!black] at (-6.0,-1.8) {Tier 2};
  \node[font=\footnotesize, draw=green!50!black, fill=white, rounded corners=3pt,
        minimum width=3.0cm, minimum height=0.65cm] at (-2.5,-1.8) {FastAPI (Python 3.12)};
  \node[font=\footnotesize, draw=green!40!black, fill=green!5, rounded corners=3pt,
        minimum width=3.0cm, minimum height=0.65cm] at (1.2,-1.8)  {Workers asynchrones};
  \node[font=\footnotesize, draw=gray!50, fill=gray!6, rounded corners=3pt,
        minimum width=2.6cm, minimum height=0.65cm] at (4.7,-1.8)  {Pydantic v2};
  \node[font=\tiny\itshape, text=gray!70] at (0,-1.05) {Couche API};

  % ── Tier 3 : Données ──
  \fill[purple!8,  rounded corners=6pt] (-7.2,-4.3) rectangle (7.2,-2.9);
  \draw[purple!50, rounded corners=6pt] (-7.2,-4.3) rectangle (7.2,-2.9);
  \node[font=\small\bfseries, text=purple!70] at (-6.0,-3.6) {Tier 3};
  \node[font=\footnotesize, draw=purple!50, fill=white, rounded corners=3pt,
        minimum width=3.0cm, minimum height=0.65cm] at (-2.5,-3.6) {PostgreSQL 17 + RLS};
  \node[font=\footnotesize, draw=purple!45, fill=purple!5, rounded corners=3pt,
        minimum width=3.0cm, minimum height=0.65cm] at (1.2,-3.6) {Supabase Auth (JWT)};
  \node[font=\footnotesize, draw=purple!40, fill=purple!5, rounded corners=3pt,
        minimum width=2.6cm, minimum height=0.65cm] at (4.7,-3.6) {Storage (S3)};
  \node[font=\tiny\itshape, text=gray!70] at (0,-2.85) {Couche Données};

  % ── Tier 4 : Infrastructure ──
  \fill[gray!10, rounded corners=6pt] (-7.2,-6.1) rectangle (7.2,-4.7);
  \draw[gray!45, rounded corners=6pt] (-7.2,-6.1) rectangle (7.2,-4.7);
  \node[font=\small\bfseries, text=gray!65] at (-6.0,-5.4) {Tier 4};
  \node[font=\footnotesize, draw=gray!45, fill=white, rounded corners=3pt,
        minimum width=2.4cm, minimum height=0.65cm] at (-4.0,-5.4) {Docker Compose};
  \node[font=\footnotesize, draw=gray!45, fill=white, rounded corners=3pt,
        minimum width=2.4cm, minimum height=0.65cm] at (-1.2,-5.4) {Let's Encrypt TLS};
  \node[font=\footnotesize, draw=gray!45, fill=white, rounded corners=3pt,
        minimum width=2.4cm, minimum height=0.65cm] at (1.8,-5.4)  {Sentry + Kuma};
  \node[font=\footnotesize, draw=gray!45, fill=white, rounded corners=3pt,
        minimum width=2.4cm, minimum height=0.65cm] at (4.8,-5.4)  {Backblaze B2};
  \node[font=\tiny\itshape, text=gray!70] at (0,-4.65) {Couche Infrastructure};

  % ── Flèches entre tiers ──
  \draw[->, thick, gray!40] (0,-0.7) -- (0,-1.1);
  \draw[->, thick, gray!40] (0,-2.5) -- (0,-2.9);
  \draw[->, thick, gray!40] (0,-4.3) -- (0,-4.7);
\end{tikzpicture}
\caption{Architecture globale de VitaChain --- diagramme en couches N-tiers}
\label{fig:2.8}
\end{figure}

\paragraph{Couche Présentation (Tier 1).}
\textbf{Frontend Next.js~15} : interface principale (agriculteurs, restaurateurs,
admin). Composants serveur (RSC) pour les lectures d'état auth. Middleware Supabase
SSR pour la protection des routes \code{/dashboard}. \textbf{Application SecondServe
(Vite/React)} : SPA indépendante pour les citoyens.

\paragraph{Couche API (Tier 2).}
\textbf{FastAPI (Python 3.12)} : API REST versionnée (\code{/api/v1/}). Validation
Pydantic~v2. Authentification Bearer JWT. Routeurs modulaires par domaine.
\textbf{Workers asynchrones} : processus Python dédiés pour les tâches longues
(diagnostics IA, emails, détection hors ligne, expiration d'annonces).

\paragraph{Couche Données (Tier 3).}
\textbf{PostgreSQL~17 (Supabase)} : base de données relationnelle avec RLS,
fonctions SECURITY DEFINER, LISTEN/NOTIFY. \textbf{Supabase Storage} : stockage
objet pour photos d'annonces et documents KYC. \textbf{Supabase Auth} : gestion
des utilisateurs, sessions JWT, Custom Access Token Hook.

\paragraph{Couche Infrastructure (Tier 4).}
\textbf{NGINX~1.27} : reverse proxy, terminaison TLS, rate-limiting.
\textbf{Docker Compose} : orchestration des services. \textbf{Backblaze B2} :
sauvegardes nocturnes. \textbf{Sentry} : capture d'erreurs.
\textbf{Uptime Kuma} : monitoring de disponibilité.

\begin{figure}[H]
    \centering
    \includegraphics[width=1\textwidth]{deploiment.png}
    \caption{Diagramme de déploiement VitaChain — infrastructure Docker/VPS}
    \label{fig:fig1_2}
\end{figure}

\section{Conclusion}

Ce chapitre a présenté l'analyse et la conception de VitaChain. L'identification des
sous-modules a permis de délimiter clairement les responsabilités de chaque composant.
Les cas d'utilisation ont formalisé les interactions entre acteurs et système. Le
modèle physique de données, avec ses 12 tables, ses fonctions SECURITY DEFINER et ses
vues d'agrégation, constitue le socle de sécurité et de performance. Les design
patterns adoptés garantissent la modularité et la maintenabilité de la solution. Le
chapitre suivant détaille la phase de réalisation.

% =============================================================
% CHAPITRE 3 — RÉALISATION
% =============================================================
\chapter{Réalisation}

\section{Introduction}

Ce chapitre présente la réalisation concrète du projet VitaChain. Nous y décrivons la
méthodologie de développement adoptée, les sprints réalisés, les technologies et outils
utilisés, ainsi qu'une présentation des interfaces de chaque module.

\section{Planning réel}

\subsection{Méthodologie adoptée}

Le développement de VitaChain a suivi une méthodologie \textbf{Scrum} adaptée au
contexte d'un projet académique. Les principes appliqués sont :

\begin{itemize}
  \item \textbf{Développement itératif :} 10 sprints de 2 semaines.
  \item \textbf{Story-driven development :} chaque fonctionnalité est définie par une
        user story formalisée dans \code{docs/stories/} avec critères d'acceptation et
        définition of done.
  \item \textbf{Intégration continue :} pipeline GitHub Actions en 5 jobs (frontend,
        backend, db, infra, détection de fuites de secrets) à chaque push.
  \item \textbf{Tests automatisés :} pytest (backend), pgTAP (RLS), tests E2E (flux
        KYC et multi-parcelles).
  \item \textbf{Pre-commit hooks :} ruff, ESLint, shellcheck, vérification AUTH-05
        avant chaque commit.
\end{itemize}
\begin{figure}[H]
    \centering
    \includegraphics[width=0.8\textwidth]{scrum_process.png}
    \caption{Diagramme de processus Scrum adapté --- VitaChain}
    \label{fig:3.1_scrum}
\end{figure}

\subsection{Tableau des sprints réalisés}

\begin{table}[H]
  \centering
  \caption{Sprints réalisés --- VitaChain}
  \label{tab:sprints}
  \small
  \begin{tabularx}{\linewidth}{lccX}
    \toprule
    \textbf{Sprint} & \textbf{Durée} & \textbf{Stories} & \textbf{Résultat} \\
    \midrule
    Sprint 0 --- Infrastructure & 2 sem. & INF-01 à INF-08
      & VPS, Supabase, CI/CD, Docker, NGINX, HTTPS, backups, monitoring. \\
    Sprint 1 --- Authentification & 2 sem. & AUTH-01 à AUTH-07
      & Inscription, connexion, JWT, RLS, KYC, matrice de droits. \\
    Sprint 2 --- Katara Parcelles & 2 sem. & KAT-01 à KAT-03
      & Registre parcelles, appairage ESP32, ingestion télémétrie. \\
    Sprint 3 --- Katara Tableaux & 2 sem. & KAT-04, KAT-05, KAT-08
      & Graphiques, historique, multi-parcelles. \\
    Sprint 4 --- Katara Services & 2 sem. & KAT-06, KAT-07, KAT-09--14
      & Alertes email, diagnostics IA, météo, NDVI, détecteur offline. \\
    Sprint 5 --- FarMarket Annonces & 2 sem. & FAR-01, FAR-02, FAR-07, FAR-08
      & Création annonces, photos, catalogue, expiration. \\
    Sprint 6 --- FarMarket Commandes & 2 sem. & FAR-03, FAR-04, FAR-05
      & Commandes multi-articles, anonymisation, suivi statut. \\
    Sprint 7 --- FarMarket Avancé & 2 sem. & FAR-06, FAR-09, FAR-10
      & Paiement COD, évaluations, tableau de bord admin. \\
    Sprint 8 --- SecondServe & 2 sem. & SEC-01 à SEC-04
      & Offres surplus, commandes citoyens, SSO magic-link. \\
    Sprint 9 --- Administration & 2 sem. & ADM-01 à ADM-04
      & File KYC, gestion utilisateurs, modération FarMarket. \\
    Sprint 10 --- Tests \& Livraison & 2 sem. & ---
      & Tests de charge (Locust), corrections, documentation. \\
    \bottomrule
  \end{tabularx}
\end{table}

\subsection{Répartition des tâches et charges}

\begin{table}[H]
  \centering
  \caption{Répartition des charges par composant}
  \label{tab:charges}
  \begin{tabularx}{\linewidth}{Xcc}
    \toprule
    \textbf{Composant} & \textbf{Charge estimée (h)} & \textbf{\%} \\
    \midrule
    Infrastructure \& DevOps & 60  & 13 \\
    Base de données           & 80  & 17 \\
    Backend FastAPI           & 120 & 26 \\
    Frontend Next.js          & 100 & 22 \\
    Application SecondServe   & 40  &  9 \\
    IoT / ESP32               & 20  &  4 \\
    Tests \& Documentation    & 40  &  9 \\
    \midrule
    \textbf{Total}            & \textbf{460} & \textbf{100} \\
    \bottomrule
  \end{tabularx}
\end{table}

\begin{figure}[H]
\centering
\begin{tikzpicture}
\pie[
  radius       = 3.6,
  text         = legend,
  after number = \%,
  color = {
    blue!80,
    blue!50,
    green!65!black,
    green!45!black,
    orange!70,
    yellow!55!orange,
    gray!55
  },
  explode = {0, 0, 0.08, 0, 0, 0, 0}
]{
  13/{Infra \& DevOps},
  17/{Base de données},
  26/{Backend FastAPI},
  22/{Frontend Next.js},
  9/{SecondServe},
  4/{IoT / ESP32},
  9/{Tests \& Docs}
}
\end{tikzpicture}
\caption{Répartition des charges de développement par composant (total~$\approx$~460h)}
\label{fig:3.2}
\end{figure}

\section{Technologies utilisées}

\subsection{Langages et Frameworks}

\begin{table}[H]
  \centering
  \caption{Langages et frameworks --- VitaChain}
  \label{tab:langages}
  \small
  \begin{tabularx}{\linewidth}{llX}
    \toprule
    \textbf{Technologie} & \textbf{Version} & \textbf{Usage} \\
    \midrule
    Python      & 3.12   & Langage du backend FastAPI et des workers. \\
    FastAPI     & 0.115  & Framework API REST asynchrone, validation Pydantic~v2. \\
    TypeScript  & 5.x    & Langage frontend (Next.js et SecondServe). \\
    Next.js     & 15.1.6 & Framework React SSR/RSC pour le frontend principal. \\
    React       & 19.0.0 & Bibliothèque UI (Next.js et SecondServe). \\
    Vite        & 6.2.0  & Build tool pour l'application SecondServe (SPA). \\
    React Router& v7     & Routage SPA dans SecondServe. \\
    SQL         & PG 17  & Requêtes et définition du schéma PostgreSQL. \\
    \bottomrule
  \end{tabularx}
\end{table}

\begin{figure}[H]
    \centering
    \includegraphics[width=0.8\textwidth]{img0.jpeg}
    \caption{Logos des principaux langages et frameworks utilisés}
    \label{fig:3.1_scrum}
\end{figure}

\subsection{Logiciels et outils de développement}

\begin{table}[H]
  \centering
  \caption{Outils de développement --- VitaChain}
  \label{tab:outils}
  \small
  \begin{tabularx}{\linewidth}{lX}
    \toprule
    \textbf{Outil} & \textbf{Usage} \\
    \midrule
    VS Code          & Éditeur principal (extensions Python, TypeScript, Tailwind). \\
    Git / GitHub     & Contrôle de version, collaboration, CI/CD. \\
    GitHub Actions   & Pipeline CI/CD automatisé (5 jobs). \\
    Docker / Compose & Conteneurisation et orchestration des services. \\
    Postman          & Test manuel des endpoints API REST. \\
    Locust           & Tests de charge sur l'endpoint d'ingestion Katara. \\
    DBeaver          & Administration et inspection de la base de données. \\
    Leaflet          & Cartes interactives pour la sélection des parcelles. \\
    Arduino IDE      & Développement du firmware ESP32. \\
    Ruff             & Linter et formateur Python ultra-rapide. \\
    ESLint           & Linter TypeScript/JavaScript. \\
    pytest           & Framework de tests unitaires Python. \\
    \bottomrule
  \end{tabularx}
\end{table}

\begin{figure}[H]
    \centering
    \includegraphics[width=0.8\textwidth]{img1.jpeg}
    \caption{Logos des principaux outils de développement}
    \label{fig:3.1_scrum}
\end{figure}

\subsection{Serveurs et services Cloud}

\begin{table}[H]
  \centering
  \caption{Services Cloud et infrastructure --- VitaChain}
  \label{tab:cloud}
  \small
  \begin{tabularx}{\linewidth}{lX}
    \toprule
    \textbf{Service} & \textbf{Rôle} \\
    \midrule
    Supabase
      & Backend-as-a-Service : PostgreSQL~17, Auth (JWT), Storage, Realtime. \\
    VPS Linux
      & Hébergement de production : Docker Compose (frontend, backend,
        nginx, certbot, kuma, backup). \\
    NGINX 1.27
      & Reverse proxy, terminaison TLS, rate-limiting, compression gzip. \\
    Let's Encrypt / Certbot
      & Certificats TLS/HTTPS avec renouvellement automatique. \\
    Backblaze B2
      & Stockage objet pour les sauvegardes nocturnes (pg\_dump via rclone). \\
    Sentry
      & Capture d'erreurs et tracking des releases (frontend + backend). \\
    Uptime Kuma
      & Monitoring de disponibilité, alertes Telegram/email. \\
    OpenWeatherMap API
      & Prévisions météo pour les parcelles Katara. \\
    Sentinel-2 (Copernicus)
      & Images satellites NDVI pour l'analyse végétale des parcelles. \\
    Google Gemini API
      & Intelligence artificielle pour les diagnostics agronomiques. \\
    Brevo (Sendinblue)
      & Service d'envoi d'emails transactionnels. \\
    \bottomrule
  \end{tabularx}
\end{table}

\begin{figure}[H]
    \centering
    \includegraphics[width=\textwidth]{archi_cloud.png}
    \caption{Architecture Cloud --- diagramme des services externes intégrés}
    \label{fig:3.5_cloud}
\end{figure}

\section{Interfaces de l'application}

\subsection{Module Authentification}

Le module d'authentification est le point d'entrée de VitaChain. Il propose une
interface claire et sécurisée pour l'inscription, la connexion et la vérification
d'identité professionnelle.

\paragraph{Page d'inscription.}
La page présente un formulaire en deux panneaux : à gauche, un aperçu de l'écosystème
VitaChain ; à droite, le formulaire de création de compte avec sélection du rôle
(Agriculteur / Restaurateur / Citoyen).

\begin{figure}[H]
    \centering
    \includegraphics[width=\textwidth]{kat1.png}
    \caption{Interface — Page d’inscription VitaChain avec sélecteur de rôle}
    \label{fig:3.5_inscription}
\end{figure}

\paragraph{Page de connexion.}

\begin{figure}[H]
    \centering
    \includegraphics[width=\textwidth]{kat2.png}
    \caption{Interface — Page de connexion VitaChain}
    \label{fig:3.6_connexion}
\end{figure}
\paragraph{Page d'onboarding KYC.}
Après inscription, les professionnels sont invités à soumettre leurs documents de
vérification. L'interface guide l'utilisateur étape par étape : sélection du type de
document, téléchargement du fichier, confirmation de soumission.

\begin{figure}[H]
    \centering
    \includegraphics[width=\textwidth]{kat3.png}
    \caption{Interface — Formulaire de soumission KYC avec upload de documents (1/4)}
    \label{fig:3.7_kyc1}
\end{figure}
\begin{figure}[H]
    \centering
    \includegraphics[width=\textwidth]{kat4.png}
    \caption{Interface — Formulaire de soumission KYC (2/4)}
    \label{fig:3.7_kyc2}
\end{figure}
\begin{figure}[H]
    \centering
    \includegraphics[width=\textwidth]{kat5.png}
    \caption{Interface — Formulaire de soumission KYC (3/4)}
    \label{fig:3.7_kyc3}
\end{figure}
\begin{figure}[H]
    \centering
    \includegraphics[width=\textwidth]{kat6.png}
    \caption{Interface — Formulaire de soumission KYC (4/4)}
    \label{fig:3.7_kyc4}
\end{figure}

\subsection{Module Katara --- Surveillance IoT}

\paragraph{Tableau de bord Katara (vue multi-parcelles).}
La page d'accueil du module Katara présente une vue synthétique de toutes les parcelles
avec les indicateurs clés : dernières mesures de sol, statut des dispositifs, nombre
d'alertes actives et résumé météo.

\begin{figure}[H]
    \centering
    \includegraphics[width=\textwidth]{kat7.png}
    \caption{Interface — Tableau de bord Katara multi-parcelles avec KPI strip (1/2)}
    \label{fig:3.9_katara1}
\end{figure}
\begin{figure}[H]
    \centering
    \includegraphics[width=\textwidth]{kat8.png}
    \caption{Interface — Tableau de bord Katara multi-parcelles avec KPI strip (2/2)}
    \label{fig:3.9_katara2}
\end{figure}

\paragraph{Création d'une parcelle.}
Le formulaire intègre une carte Leaflet permettant à l'agriculteur de dessiner le
contour géographique exact de son champ (polygone GeoJSON).

\begin{figure}[H]
    \centering
    \includegraphics[width=\textwidth]{kat9.png}
    \caption{Interface --- Formulaire de création de parcelle avec carte Leaflet interactive}
    \label{fig:3.10_parcelle}
\end{figure}

\paragraph{Télémétrie en temps réel.}
Chaque parcelle dispose d'une vue détaillée avec des onglets. L'onglet télémétrie
affiche des graphiques Sparkline pour les 4 métriques de sol sur les périodes 24h, 7j
et 30j.

\begin{figure}[H]
    \centering
    \includegraphics[width=\textwidth]{kat10.png}
    \caption{Interface --- Graphiques de télémétrie temps réel pour une parcelle (4 métriques)}
    \label{fig:3.11_telemetrie}
\end{figure}

\paragraph{Gestion des seuils d'alerte.}
L'agriculteur configure des seuils min/max pour chaque métrique. Le dépassement
déclenche automatiquement un email d'alerte via le worker Katara.

\begin{figure}[H]
    \centering
    \includegraphics[width=\textwidth]{kat11.png}
    \caption{Interface --- Configuration des seuils d'alerte par métrique}
    \label{fig:3.12_seuils}
\end{figure}

\paragraph{Diagnostic IA.}
Le système agrège les données télémétriques, météorologiques et satellitaires, puis
génère un résumé agronomique en langage naturel via Gemini.

\begin{figure}[H]
    \centering
    \includegraphics[width=\textwidth]{kat12.png}
    \caption{Interface --- Résultat d'un diagnostic agronomique IA (Gemini)}
    \label{fig:3.13_diagnostic}
\end{figure}

\paragraph{Vue météo et satellite NDVI.}

\begin{figure}[H]
    \centering
    \includegraphics[width=\textwidth]{kat13.png}
    \caption{Interface --- Prévisions météo OpenWeatherMap pour une parcelle}
    \label{fig:3.14_meteo}
\end{figure}

\begin{figure}[H]
    \centering
    \includegraphics[width=\textwidth]{kat14.png}
    \caption{Interface --- Image NDVI Sentinel-2 pour une parcelle agricole}
    \label{fig:3.15_ndvi}
\end{figure}

\paragraph{Gestion des dispositifs ESP32.}

\begin{figure}[H]
    \centering
    \includegraphics[width=\textwidth]{kat15.png}
    \caption{Interface --- Liste des dispositifs IoT avec statut (ACTIVE/OFFLINE)}
    \label{fig:3.16_dispositifs}
\end{figure}

\subsection{Module FarMarket --- Marketplace Agricole}

\paragraph{Tableau de bord agriculteur --- Mes annonces.}

\begin{figure}[H]
    \centering
    \includegraphics[width=\textwidth]{farm0.png}
    \caption{Interface --- Tableau de bord agriculteur FarMarket avec liste des annonces}
    \label{fig:3.16_dispositifs}
\end{figure}

\paragraph{Création d'une annonce avec photos.}
Le formulaire permet d'uploader jusqu'à 5 photos du produit. Les photos sont stockées
dans le bucket Supabase \code{farmarket-photos} avec politiques RLS.

\begin{figure}[H]
    \centering
    \includegraphics[width=\textwidth]{farm3.png}
    \caption{Interface --- Formulaire de création d'annonce agricole avec galerie photos}
    \label{fig:3.16_dispositifs}
\end{figure}


\paragraph{Catalogue restaurateur.}
Le restaurateur accède à un catalogue paginé d'annonces actives avec filtres par
région (12 régions du Maroc), type de produit et fourchette de prix. Une recherche
textuelle floue (pg\_trgm) est disponible.

\begin{figure}[H]
    \centering
    \includegraphics[width=\textwidth]{farm4.png}
    \caption{Interface --- Catalogue FarMarket avec filtres et cartes d'annonces}
    \label{fig:3.16_dispositifs}
\end{figure}


\paragraph{Détail d'une annonce.}

\begin{figure}[H]
    \centering
    \includegraphics[width=\textwidth]{farm5.png}
    \caption{Interface --- Page détail d'une annonce avec profil agriculteur, note moyenne et avis}
    \label{fig:3.16_dispositifs}
\end{figure}

\paragraph{Panier et processus de commande.}
Le restaurateur ajoute des articles au panier depuis plusieurs agriculteurs différents.
Le système calcule automatiquement les frais logistiques selon la région de livraison.

\begin{figure}[H]
    \centering
    \includegraphics[width=\textwidth]{farm6.png}
    \caption{Interface --- Panier restaurateur avec résumé de commande et frais logistiques}
    \label{fig:3.16_dispositifs}
\end{figure}


\paragraph{Suivi de commande --- vue restaurateur.}





\paragraph{Articles entrants --- vue agriculteur (anonymisée).}
L'agriculteur voit ses articles entrants via la vue \code{v\_farmer\_incoming\_items}.
Le nom du restaurateur est masqué, remplacé par un \code{resto\_handle} opaque
(sha256).

\paragraph{Confirmation de paiement COD.}
\begin{figure}[H]
    \centering
    \includegraphics[width=\textwidth]{farm7.png}
    \caption{Interface --- Confirmation de paiement en espèces (COD) par le restaurateur}
    \label{fig:3.16_dispositifs}
\end{figure}

\paragraph{Laisser une évaluation.}

\subsection{Module SecondServe --- Anti-gaspillage}

SecondServe est une application distincte (Vite/React), partagée sur le même projet
Supabase, dédiée à la redistribution des surplus alimentaires des restaurants aux
citoyens.

\begin{figure}[H]
    \centering
    \includegraphics[width=\textwidth]{ss1.png}
    \caption{Interface --- Page d'accueil SecondServe avec liste des offres surplus disponibles}
    \label{fig:3.16_dispositifs}
\end{figure}


\paragraph{Catalogue des offres.}
Le citoyen filtre les offres par ville et par type de cuisine. Chaque carte d'offre
affiche le nom du plat, le prix réduit, la quantité restante et l'heure limite de
réservation.

\begin{figure}[H]
    \centering
    \includegraphics[width=\textwidth]{ss3.png}
    \caption{Interface --- Catalogue des offres surplus avec filtres par ville}
    \label{fig:3.16_dispositifs}
\end{figure}
\begin{figure}[H]
    \centering
    \includegraphics[width=\textwidth]{ss4.png}
    \caption{Interface --- Tableau de bord citoyen SecondServe avec historique des commandes}
    \label{fig:3.16_dispositifs}
\end{figure}

\begin{figure}[H]
    \centering
    \includegraphics[width=\textwidth]{ss5.png}
    \caption{Interface --- Tableau de bord restaurateur SecondServe (gestion des offres)}
    \label{fig:3.16_dispositifs}
\end{figure}

\begin{figure}[H]
    \centering
    \includegraphics[width=\textwidth]{ss6.png}
    \caption{Interface --- Reçu de commande SecondServe}
    \label{fig:3.16_dispositifs}
\end{figure}


\subsection{Module Administration}

\paragraph{File d'attente KYC.}
L'administrateur voit la liste des demandes en attente avec accès aux documents soumis.
Il peut approuver ou rejeter chaque demande. L'action déclenche la mise à jour du
statut et l'envoi d'un email.

\begin{figure}[H]
    \centering
    \includegraphics[width=\textwidth]{AD1.png}
    \caption{Interface --- File d'attente KYC avec téléchargement de documents et boutons d'action}
    \label{fig:3.16_dispositifs}
\end{figure}


\paragraph{Gestion des utilisateurs.}

\begin{figure}[H]
    \centering
    \includegraphics[width=\textwidth]{AD2.png}
    \caption{Interface --- Console de gestion des utilisateurs avec filtres et actions}
    \label{fig:3.16_dispositifs}
\end{figure}


\paragraph{Tableau de bord FarMarket (administration).}

\begin{figure}[H]
    \centering
    \includegraphics[width=\textwidth]{AD3.png}
    \caption{Interface --- Tableau de bord admin FarMarket avec statistiques globales}
    \label{fig:3.16_dispositifs}
\end{figure}
\begin{figure}[H]
    \centering
    \includegraphics[width=\textwidth]{AD4.png}
    \caption{Interface --- Tableau de bord admin SecondServe avec statistiques globales}
    \label{fig:3.16_dispositifs}
\end{figure}


\subsection{Dispositif IoT --- ESP32}

Le dispositif IoT au cœur du module Katara est un microcontrôleur \textbf{ESP32}
équipé de capteurs de sol mesurant les paramètres agronomiques et les transmettant
au backend VitaChain via HTTP.

\paragraph{Capteurs intégrés :}
capteur d'humidité du sol (capacitif) ; capteur de température du sol (DS18B20) ;
capteur de pH du sol (électrode pH + module analogique) ; capteur de conductivité
électrique (EC meter) ; mesure de niveau de batterie (diviseur résistif sur ADC).

\fig{fig3_34_esp32_schema}{Schéma du dispositif IoT --- ESP32 avec capteurs de sol connectés}{3.34}

\paragraph{Protocole de communication.}
Le firmware envoie une requête HTTP POST toutes les $N$ minutes (configurable) vers
l'endpoint \api{POST /api/v1/katara/ingest}.

\begin{lstlisting}[style=json,caption={Payload télémétrique ESP32},label=lst:payload]
{
  "soil_moisture":     42.5,
  "soil_temperature":  21.3,
  "soil_ph":            6.8,
  "soil_conductivity":  1.2,
  "battery_level":     78.0,
  "recorded_at": "2026-06-17T10:30:00Z"
}
\end{lstlisting}

\begin{lstlisting}[style=bash,caption={En-têtes HTTP d'authentification du dispositif},label=lst:headers]
X-Device-Id:      ESP-KAT-001
X-Device-Api-Key: vk_a1b2c3d4e5f6...
\end{lstlisting}

\paragraph{Flux d'appairage.}

\begin{figure}[H]
    \centering
    \includegraphics[width=0.8\textwidth]{fig3_35_appairage_seq1.png}
    \caption{Diagramme de séquence --- Flux d'appairage d'un ESP32 avec VitaChain (1/3)}
    \label{fig:3.35_appairage1}
\end{figure}
\begin{figure}[H]
    \centering
    \includegraphics[width=0.8\textwidth]{fig3_35_appairage_seq2.png}
    \caption{Diagramme de séquence --- Flux d'appairage d'un ESP32 avec VitaChain (2/3)}
    \label{fig:3.35_appairage2}
\end{figure}
\begin{figure}[H]
    \centering
    \includegraphics[width=0.8\textwidth]{fig3_35_appairage_seq3.png}
    \caption{Diagramme de séquence --- Flux d'appairage d'un ESP32 avec VitaChain (3/3)}
    \label{fig:3.35_appairage3}
\end{figure}

\fig{fig3_36_esp32_photo}{Photo du prototype ESP32 avec capteurs de sol assemblés sur breadboard}{3.36}

\section{Conclusion}

Ce chapitre a présenté la réalisation complète de VitaChain. La méthodologie Scrum a
permis de livrer 50+ stories réparties sur 10 sprints en environ 460 heures de
développement. La plateforme intègre un stack technologique moderne et couvre quatre
modules fonctionnels complets. Le dispositif ESP32 assure la collecte des données
agronomiques en temps réel avec une latence d'ingestion inférieure à 50~ms.

% =============================================================
% CHAPITRE 4 — CONCLUSION GÉNÉRALE
% =============================================================
\chapter{Conclusion Générale}

\section{Bilan}

VitaChain est une plateforme numérique agroalimentaire complète, conçue et réalisée de
A~à~Z dans le cadre de ce Projet de Fin d'Études. Le projet a atteint l'ensemble de ses
objectifs initiaux.

\paragraph{Objectifs techniques atteints.}
\begin{itemize}
  \item Architecture full-stack production-ready : FastAPI, Next.js~15, PostgreSQL~17,
        Docker, NGINX, TLS.
  \item Système d'authentification robuste avec JWT multi-rôle, Row-Level Security
        (22 politiques) et processus KYC.
  \item Module Katara opérationnel : ingestion IoT $<$ 50~ms (p50), alertes temps réel
        via LISTEN/NOTIFY, diagnostics IA, intégration OWM et Sentinel-2.
  \item Module FarMarket complet : cycle de vie des annonces, commandes
        multi-agriculteurs, anonymisation B2B, paiement COD, évaluations.
  \item Module SecondServe fonctionnel : offres surplus, commandes citoyens, SSO
        inter-application via magic-link.
  \item Administration centralisée : KYC, gestion utilisateurs, modération.
  \item Infrastructure DevOps : 50 migrations SQL, CI/CD GitHub Actions (5 jobs),
        sauvegardes nocturnes, monitoring Sentry + Uptime Kuma.
\end{itemize}

\paragraph{Défis surmontés.}
\begin{itemize}
  \item \textbf{Sécurité multi-rôle complexe :} conception et test de la matrice RLS
        22-cellules pour éviter les fuites de données inter-utilisateurs.
  \item \textbf{Performance IoT :} fonction SQL atomique SECURITY DEFINER combinant
        vérification, insertion, dénormalisation et notification en un seul appel pour
        atteindre la cible $<$ 50~ms.
  \item \textbf{Anonymisation B2B :} vue \code{v\_farmer\_incoming\_items} avec hash
        sha256 déterministe préservant l'utilité tout en protégeant l'identité du
        restaurateur.
  \item \textbf{SSO inter-application :} handoff SecondServe par OTP Supabase à usage
        unique.
\end{itemize}

\paragraph{Chiffres clés du projet.}
\begin{itemize}
  \item 50 migrations SQL ; 12 tables PostgreSQL ; 5 vues de sécurité
  \item $\approx$~60 endpoints REST API
  \item 4 modules fonctionnels livrés
  \item $\approx$~460 heures de développement sur 10 sprints
\end{itemize}

\section{Perspectives}

\paragraph{Court terme.}
\begin{itemize}
  \item \textbf{Module BotaBa9a :} finaliser l'implémentation du module de surveillance
        de chaîne du froid, dont l'architecture est déjà définie.
  \item \textbf{Notifications push :} ajouter des notifications push web (Web Push API)
        en complément des emails pour les alertes Katara et mises à jour FarMarket.
  \item \textbf{Application mobile :} développer une application React Native pour
        améliorer l'accessibilité aux agriculteurs en zone rurale.
\end{itemize}

\paragraph{Moyen terme.}
\begin{itemize}
  \item \textbf{Paiement en ligne :} intégrer une vraie passerelle de paiement
        (CMI, PayZone Maroc).
  \item \textbf{Intelligence artificielle avancée :} enrichir les diagnostics Katara
        avec des modèles ML spécialisés (prédiction de rendement, recommandation
        d'irrigation).
  \item \textbf{Traçabilité :} explorer l'ajout d'une couche de traçabilité sur la
        chaîne d'approvisionnement FarMarket.
\end{itemize}

\paragraph{Long terme.}
\begin{itemize}
  \item \textbf{Marketplace de données agronomiques :} permettre aux agriculteurs de
        monétiser leurs données anonymisées auprès d'instituts de recherche.
  \item \textbf{Intégration API nationale :} connecter VitaChain aux systèmes
        d'information du Ministère de l'Agriculture marocain (MAPMDREF).
  \item \textbf{Réseau de capteurs mutualisé :} déployer des nœuds IoT communautaires
        permettant à plusieurs agriculteurs d'une même zone de partager
        l'infrastructure de capteurs.
\end{itemize}

VitaChain démontre la faisabilité technique et la pertinence d'une plateforme numérique
intégrée pour le secteur agroalimentaire marocain. La combinaison de l'IoT, du
marketplace B2B et de l'anti-gaspillage dans un seul écosystème constitue une
proposition de valeur différenciante et un point de départ solide pour un déploiement
à grande échelle.

% =============================================================
% ANNEXES
% =============================================================
\appendix

\chapter{Configuration du dispositif ESP32}
\label{annexe:a}

\section{Schéma de câblage des capteurs}

\fig{figA_1_cablage}{Schéma de câblage complet ESP32 --- capteurs de sol, alimentation et connectivité Wi-Fi}{A.1}

\section{Format du payload télémétrique}

Le dispositif ESP32 envoie un payload JSON à intervalle configurable (défaut :
5~minutes) :

\begin{lstlisting}[style=json,caption={Payload JSON complet envoyé par l'ESP32},label=lst:payloadfull]
{
  "soil_moisture":      45.2,
  "soil_temperature":   22.1,
  "soil_ph":             6.75,
  "soil_conductivity":   1.35,
  "battery_level":      82.0,
  "recorded_at": "2026-06-17T10:30:00.000Z"
}
\end{lstlisting}

En-têtes HTTP requis :
\begin{lstlisting}[style=bash,label=lst:headersannex]
Content-Type:     application/json
X-Device-Id:      ESP-KAT-001
X-Device-Api-Key: vk_<32_hex_chars>
\end{lstlisting}

Réponse attendue en succès :
\begin{lstlisting}[style=json,label=lst:response]
{"status": "ok", "recorded_at": "2026-06-17T10:30:00.000Z"}
\end{lstlisting}

\section{Procédure d'appairage d'un nouveau dispositif}

\begin{enumerate}
  \item L'agriculteur crée une parcelle dans l'interface VitaChain.
  \item Il accède à l'onglet \og Dispositifs \fg{} et clique sur
        \og Associer un capteur \fg{}.
  \item Il saisit l'identifiant du dispositif imprimé sur le boîtier ESP32
        (format \code{ESP-KAT-NNN}).
  \item Le backend génère la clé d'API et l'affiche \textbf{une seule fois}.
  \item L'agriculteur saisit la clé dans le fichier \code{devices.json} de l'ESP32.
  \item L'ESP32 envoie sa première mesure ; le dispositif passe automatiquement en
        statut ACTIVE.
\end{enumerate}

\section{Codes d'erreur de l'endpoint d'ingestion}

\begin{table}[H]
  \centering
  \caption{Codes d'erreur de \texttt{POST /api/v1/katara/ingest}}
  \label{tab:erreurs}
  \begin{tabularx}{\linewidth}{clX}
    \toprule
    \textbf{Code HTTP} & \textbf{Erreur} & \textbf{Cause} \\
    \midrule
    200 & OK                        & Mesure enregistrée avec succès. \\
    401 & invalid\_device\_credentials & Clé d'API incorrecte ou dispositif
                                         inexistant. \\
    409 & duplicate\_timestamp       & Mesure déjà reçue pour ce
                                       (device\_id, recorded\_at). \\
    422 & validation\_error          & Payload JSON invalide ou valeurs hors
                                       plage. \\
    429 & rate\_limited              & Trop de requêtes (limite NGINX :
                                       30~req/s). \\
    503 & service\_unavailable       & Base de données Supabase inaccessible. \\
    \bottomrule
  \end{tabularx}
\end{table}

\chapter{Matrice des droits d'accès (RLS)}
\label{annexe:b}

\section{Matrice rôle \texorpdfstring{$\times$}{x} table \texorpdfstring{$\times$}{x} verbe}

La matrice suivante résume les politiques Row-Level Security implémentées dans
PostgreSQL pour VitaChain.

\begin{longtable}{p{5.5cm}ccccc}
  \caption{Matrice complète RLS --- VitaChain (22 cellules vérifiées)}
  \label{tab:rls}\\
  \toprule
  \textbf{Table / Opération}
    & \textbf{FARMER} & \textbf{RESTO} & \textbf{CITIZEN}
    & \textbf{ADMIN}  & \textbf{Service} \\
  \midrule
  \endfirsthead
  \multicolumn{6}{c}{\small\itshape Suite du tableau~\ref{tab:rls}}\\
  \toprule
  \textbf{Table / Opération}
    & \textbf{FARMER} & \textbf{RESTO} & \textbf{CITIZEN}
    & \textbf{ADMIN}  & \textbf{Service} \\
  \midrule
  \endhead
  \bottomrule
  \endlastfoot

  \code{profiles} (SELECT)
    & C(own) & C(own) & C(own) & \checkmark & \checkmark \\
  \code{profiles} (UPDATE)
    & C(own) & C(own) & C(own) & \checkmark & \checkmark \\
  \code{kyc\_documents} (SELECT)
    & C(own) & C(own) & $\times$ & \checkmark & \checkmark \\
  \code{kyc\_documents} (INSERT)
    & \checkmark & \checkmark & $\times$ & $\times$ & \checkmark \\
  \code{m1\_katara\_parcels} (SELECT)
    & C(own) & $\times$ & $\times$ & \checkmark & \checkmark \\
  \code{m1\_katara\_parcels} (INSERT)
    & V & $\times$ & $\times$ & $\times$ & \checkmark \\
  \code{m1\_katara\_devices} (SELECT)
    & C(own) & $\times$ & $\times$ & \checkmark & \checkmark \\
  \code{m1\_katara\_telemetry} (SELECT)
    & C(own) & $\times$ & $\times$ & \checkmark & \checkmark \\
  \code{m1\_katara\_telemetry} (INSERT)
    & $\times$ & $\times$ & $\times$ & $\times$ & \checkmark \\
  \code{m2\_farmarket\_ads} (SELECT)
    & \checkmark & \checkmark & \checkmark & \checkmark & \checkmark \\
  \code{m2\_farmarket\_ads} (INSERT)
    & V & $\times$ & $\times$ & $\times$ & \checkmark \\
  \code{m2\_farmarket\_orders} (SELECT)
    & $\times$ & C(own) & $\times$ & \checkmark & \checkmark \\
  \code{m2\_farmarket\_orders} (INSERT)
    & $\times$ & V & $\times$ & $\times$ & \checkmark \\
  \code{m2\_farmarket\_order\_items} (SELECT)
    & C(anon.) & C(own) & $\times$ & \checkmark & \checkmark \\
  \code{m2\_farmarket\_farmer\_ratings} (SELECT)
    & \checkmark & \checkmark & \checkmark & \checkmark & \checkmark \\
  \code{m2\_farmarket\_farmer\_ratings} (INSERT)
    & $\times$ & D & $\times$ & $\times$ & \checkmark \\
\end{longtable}

\vspace{6pt}
\noindent
\textbf{Légende :}
\begin{itemize}
  \item \checkmark{} : accès autorisé sans condition supplémentaire.
  \item $\times$ : accès interdit.
  \item C(own) : accès conditionnel --- uniquement ses propres données
        (\code{auth.uid() = user\_id}).
  \item C(anon.) : accès via la vue anonymisée
        (\code{v\_farmer\_incoming\_items}).
  \item V : autorisé uniquement si \code{verification\_status = 'VERIFIED'}
        (WITH CHECK).
  \item D : autorisé uniquement si au moins 1 article DELIVERED de cet
        agriculteur (WITH CHECK).
\end{itemize}

\section{Fonctions SECURITY DEFINER}

\begin{table}[H]
  \centering
  \caption{Fonctions PostgreSQL SECURITY DEFINER --- VitaChain}
  \label{tab:secdef}
  \small
  \begin{tabularx}{\linewidth}{llX}
    \toprule
    \textbf{Fonction} & \textbf{Signature} & \textbf{Description} \\
    \midrule
    \code{has\_role(role)}
      & \code{(user\_role)} $\rightarrow$ \code{bool}
      & Vérifie le rôle JWT sans récursion RLS. \\
    \code{is\_admin()}
      & \code{()} $\rightarrow$ \code{bool}
      & Raccourci pour \code{has\_role('ADMIN')}. \\
    \code{verify\_device\_api\_key}
      & \code{(text, text)} $\rightarrow$ \code{bool}
      & Comparaison bcrypt en temps constant. \\
    \code{m1\_katara\_ingest}
      & \code{(...)} $\rightarrow$ \code{void}
      & Ingestion atomique : vérif + INSERT + UPDATE + NOTIFY. \\
    \code{ss\_place\_order}
      & \code{(...)} $\rightarrow$ \code{uuid}
      & Commande SecondServe avec contrôle de stock atomique. \\
    \bottomrule
  \end{tabularx}
\end{table}

\end{document}
